%%%%%%%%%%%%%%%%%%%%%%%%%%%%%%%%%%%%%%%%%%%%%%%%%%%%%%%%%%%%%%%%%%%%%%%%%%%%%%%%
%%  ENDABGABE PORTFOLIO  (15.03.2027)  – Hauptdatei
%%  Aufbau nach pp-portfolio-geruest-2026-2027.tex:
%%    Titelseite – Inhaltsverzeichnis – Einleitung – Hauptteil (6 Pflichteinlagen
%%    + freiwillige) – Gesamtreflexion – Literaturverzeichnis – Anhänge
%%  Texte schreibst du in kapitel/ und einlagen/, Angaben in einstellungen.tex.
%%  Overleaf: Menü -> Main document = diese Datei (für Endabgabe).
%%%%%%%%%%%%%%%%%%%%%%%%%%%%%%%%%%%%%%%%%%%%%%%%%%%%%%%%%%%%%%%%%%%%%%%%%%%%%%%%
%%%%%%%%%%%%%%%%%%%%%%%%%%%%%%%%%%%%%%%%%%%%%%%%%%%%%%%%%%%%%%%%%%%%%%%%%%%%%%%%
%%  EINSTELLUNGEN – die einzige Datei, in der du deine Angaben einträgst.
%%  Wird von ALLEN Hauptdateien (Endabgabe + Zwischenabgaben) eingelesen.
%%  Entspricht dem Block "Die Angaben" der offiziellen Gerüste
%%  pp-portfolio-geruest-2026-2027.tex / pp-portfolio-geruest-einlage-2026-2027.tex
%%%%%%%%%%%%%%%%%%%%%%%%%%%%%%%%%%%%%%%%%%%%%%%%%%%%%%%%%%%%%%%%%%%%%%%%%%%%%%%%

%%%%%%%%%%%%%%%%%%%%%%%%%%%%%%%%%%%%% Die Angaben %%%%%%%%%%%%%%%%%%%%%%%%%%%%%%%
\newcommand{\Semester}{WH 2026-2027}
\newcommand{\Projekt}{Dynamik Keltischer Wackelsteine}   % Titel laut Themenliste WS 26/27
%% Schwerpunkt abgeleitet aus der Aufgabenverteilung im Arbeitsplan (1.4 Bewegungsgleichungen,
%% 1.6 Simulation Umkehrdynamik, 2.1.3/2.1.4 Firmware + Kalibrierung, 5.1 Datenpipeline) – prüfen!
\newcommand{\Schwerpunkt}{Bewegungsgleichungen und numerische Simulation der Spin-Umkehr}
\newcommand{\SchwerpunktKurz}{Simulation der Spin-Umkehr}        % kurze Form für die Kopfzeile
\newcommand{\IhrName}{Kristijan Manevski}                         % Vorname Nachname
\newcommand{\Matrikelnummer}{\platzhalter{Matrikelnummer}}
\newcommand{\Studiengang}{\platzhalter{Studiengang + Semester}}   % z.B. Physik B.Sc., 5. Semester
\newcommand{\emailAdresse}{\platzhalter{Uni-E-Mail-Adresse (\dots @uni-ulm.de)}}
%% Abgabedatum: jede Zwischenabgabe-Hauptdatei setzt ihr eigenes \dasDatum (Moodle-Frist);
%% ohne Vorgabe gilt das Datum der Endabgabe.
\providecommand{\dasDatum}{15.\,03.\,2027}

%% Weitere Angaben, die in mehreren Einlagen gebraucht werden
\newcommand{\Partner}{Janosch Lustig}                             % für Abgrenzung der Schwerpunkte
%% laut Titelseite von Janoschs Arbeitsplan-Abgabe (28.10.2026)
\newcommand{\PartnerSchwerpunkt}{Umkehrzeit und Parameterabhängigkeit}
\newcommand{\Laborbuchlink}{https://link-zum-elektronischen-laborjournal}  % Pflicht in Einlage 3!
%% Einlage 3: ein Praktikumstag (Do 10--12 Uhr, O26/233 bzw. O26/234), ≠ Tag von Janosch.
%% Termine: 15.10., 22.10., 29.10., 05.11., 12.11., 19.11., 26.11., 03.12., 10.12., 17.12.2026,
%%          14.01., 21.01., 28.01., 04.02., 11.02.2027 (Abgabe Protokoll: 09.12.2026)
\newcommand{\Protokolldatum}{TT.MM.2026}
\newcommand{\PaperKey}{KondoNakanishi2017}      % Einlage 2: Kondo & Nakanishi, Phys. Rev. E 95, 062207
\newcommand{\Vortragsdatum}{28.\,01.\,2027}      % Einlage 5
\newcommand{\FremdVortragKey}{VORLAGE-Vortrag}  % Einlage 6: bib-Key des anderen Vortrags

%%%%%%%%%%%%%%%%%%%%%%%%%%%%%%%%%%%%%%%%%%%%%%%%%%%%%%%%%%%%%%%%%%%%%%%%%%%%%%%%
%%  MODUS-SCHALTER
%%
%%  \Abgabefalse  = Entwurf: grüne Hilfen sichtbar, Kontrollseite mit
%%                  Umfangsmessung, Platzhalter rot markiert.
%%  \Abgabetrue   = Abgabe: grüne Farbe -> schwarz (wie offizielle Zeile 15/16),
%%                  alle Hilfen + Kontrollseite + Umfangsboxen verschwinden.
%%                  Übrig gebliebene \platzhalter{} bleiben ROT sichtbar und
%%                  erzeugen eine Warnung -> dann ist noch etwas offen!
%%
%%  \Hinweisefalse blendet im Entwurf nur die grünen Hilfstexte aus
%%                  (nützlich, um den Umfang realistisch zu messen).
%%%%%%%%%%%%%%%%%%%%%%%%%%%%%%%%%%%%%%%%%%%%%%%%%%%%%%%%%%%%%%%%%%%%%%%%%%%%%%%%
\newif\ifAbgabe
\newif\ifHinweise

\Abgabefalse      % <-- für die Abgabe auf \Abgabetrue setzen
\Hinweisetrue     % <-- Hilfstexte im Entwurf ein-/ausblenden


%%%%%%%%%%%%%%%%%%%%%%%%%%%%%%%%%%%%%% TeX Front Matter %%%%%%%%%%%%%%%%%%%%%%%%%
\documentclass[pdftex,12pt,twoside,footsepline,headsepline,a4paper,final]{scrreprt}
%%%%%%%%%%%%%%%%%%%%%%%%%%%%%%%%%%%%%%%%%%%%%%%%%%%%%%%%%%%%%%%%%%%%%%%%%%%%%%%%
%%  PRÄAMBEL – gemeinsam für Endabgabe und alle Zwischenabgaben.
%%  Teil 1 ist 1:1 die Präambel der offiziellen Gerüste (nicht verändern),
%%  Teil 2 ergänzt nur Pakete für Gleichungen, Bilder, PDF-Einbindung usw.,
%%  Teil 3 enthält die Hilfsmakros dieser Vorlage.
%%  Die \documentclass-Zeile steht (unverändert offiziell) in jeder Hauptdatei.
%%%%%%%%%%%%%%%%%%%%%%%%%%%%%%%%%%%%%%%%%%%%%%%%%%%%%%%%%%%%%%%%%%%%%%%%%%%%%%%%

%%%%%%%%%%%%%%%%%%%%%%%%%%%%% Teil 1: offizielle Präambel %%%%%%%%%%%%%%%%%%%%%%
\usepackage[verbose,a4paper,bottom=25mm,top=25mm]{geometry} %showframe,
\usepackage{scrlayer-scrpage}
\usepackage[T1]{fontenc}
\usepackage[utf8]{inputenc}
\usepackage{ngerman}
\usepackage{color}
\usepackage{amsfonts}
\usepackage{amssymb}
\usepackage{amsthm}
%%%%%%%%%%%%%%%%%%%%%%%%% Biblatex (offizielle Optionen, Backend Biber) %%%%%%%%%%
%% Overleaf erkennt biber automatisch. Zitierstil alphabetic: [Sch02]
\usepackage[sorting=nyt, backend=biber, citestyle=alphabetic, style=alphabetic,
            autocite=inline, labelalpha=true, hyperref=auto, backref=true,
            language=ngerman]{biblatex}
\addbibresource{literatur.bib}

%%%%%%%%%%%%%%%%%%%%%%%%%%%%% Teil 2: Zusatzpakete %%%%%%%%%%%%%%%%%%%%%%%%%%%%%%
\usepackage{amsmath}               % \eqref, align, nummerierte Gleichungen
\usepackage{xcolor}
\definecolor{darkgreen}{rgb}{0,0.5,0}   % offizielle Hinweisfarbe
\usepackage{graphicx}              % \includegraphics
\graphicspath{{abbildungen/}}
\usepackage{booktabs}              % schöne Tabellen
\usepackage[locale=DE]{siunitx}    % \SI{3.4}{\kilo\hertz}, \num{1.2e3}
\usepackage{pdfpages}              % \includepdf für Anhänge (Paper, Laborjournal, Folien)
\usepackage{pdflscape}             % Querformat: \begin{landscape}...\end{landscape}
\usepackage{pgfgantt}              % Gantt-Diagramm für den Arbeitsplan
\usepackage{environ}
\usepackage[savepos,abspage]{zref}
\usepackage[hidelinks]{hyperref}   % klickbare Verweise; zuletzt laden

%% Zitierregel (Marti): bei Journalartikeln DOI angeben, aber keine Webadresse.
%% -> url/urldate werden bei @article automatisch unterdrückt.
\AtEveryBibitem{%
  \ifentrytype{article}{\clearfield{url}\clearfield{urlyear}%
    \clearfield{urlmonth}\clearfield{urlday}}{}}

%% Kopf- und Fußzeile wie offizielles Gerüst (\ohead wird in der Hauptdatei gesetzt)
\pagestyle{scrheadings}
\ihead{\thepage}
\ofoot{\pagemark}
\ifoot{\zubearbeiten{\IhrName}}

%%%%%%%%%%%%%%%%%%%%%%%%%%%%% Teil 3: Hilfsmakros %%%%%%%%%%%%%%%%%%%%%%%%%%%%%%%
\makeatletter

%% --- offizieller Farbschalter (Zeile 12-15 der Gerüste), automatisch per Modus
\ifAbgabe
  \Hinweisefalse
  \newcommand{\zubearbeiten}[1]{\textcolor{black}{#1}}
\else
  \newcommand{\zubearbeiten}[1]{\textcolor{darkgreen}{#1}}
\fi

%% --- Grüne Hilfstexte: verschwinden im Abgabemodus vollständig
%%     \begin{hinweise} ... \end{hinweise}   (darf itemize enthalten)
%%     \hinweis{kurzer Text}
\NewEnviron{hinweise}{\ifHinweise\par{\color{darkgreen}\small\BODY\par}\fi}
\newcommand{\hinweis}[1]{\ifHinweise\par{\color{darkgreen}\small #1\par}\fi}

%% --- Platzhalter für eigenen Text: bleibt IMMER rot sichtbar (auch im
%%     Abgabemodus), damit nichts Unfertiges abgegeben wird.
\newcounter{platzhalter}
\DeclareRobustCommand{\platzhalter}[1]{%
  \stepcounter{platzhalter}%
  \textcolor{red}{\textbf{[}\,#1\,\textbf{]}}%
  \ifAbgabe\PackageWarning{Portfolio}{Platzhalter noch vorhanden: #1}\fi}
\AtEndDocument{%
  \ifnum\value{platzhalter}>0
    \PackageWarningNoLine{Portfolio}{Es sind noch \theplatzhalter\space
      Platzhalter im Dokument!}%
  \fi
  \immediate\write\@auxout{\string\gdef\string\pf@platzhalteralt{\theplatzhalter}}}
\providecommand{\pf@platzhalteralt}{?}

%% --- Umfangsmessung
%%   \begin{umfang}{Bezeichnung}{min}{max} ... \end{umfang}
%%   misst die Länge des eingeschlossenen Textes in Seiten (über die
%%   Position auf der Seite, auf 0.05 Seiten genau) und zeigt sie im Entwurf
%%   in einem Kasten an. Achtung: Abbildungen/Tabellen zählen nach den Regeln
%%   NICHT zur Länge, werden hier aber mitgemessen -> dann abziehen.
\zref@addprop{savepos}{abspage}
\newcounter{pf@umfang}
\newcommand*{\pf@umfangliste}{}
\newcommand*{\pf@umfangeintrag}[4]{% Name, min, max, ist
  \g@addto@macro\pf@umfangliste{\pf@umfangzeile{#1}{#2}{#3}{#4}}}
\newcommand*{\pf@setfarbe}[3]{% min max ist -> \pf@farbe
  \def\pf@farbe{darkgreen}%
  \ifdim #3pt<#1pt \def\pf@farbe{red}\fi
  \ifdim #3pt>#2pt \def\pf@farbe{red}\fi}
\newcommand*{\pf@umfangzeile}[4]{%
  #1 & #2 -- #3 & \pf@setfarbe{#2}{#3}{#4}\textcolor{\pf@farbe}{\textbf{#4}}\\}
\newenvironment{umfang}[3]{%
  \stepcounter{pf@umfang}%
  \edef\pf@uid{pfumf\the\value{pf@umfang}}%
  \def\pf@uname{#1}\def\pf@umin{#2}\def\pf@umax{#3}%
  \par\zsavepos{\pf@uid a}\nobreak
}{%
  \par\zsavepos{\pf@uid e}%
  \ifAbgabe\else
    \zref@refused{\pf@uid a}\zref@refused{\pf@uid e}%
    \edef\pf@ist{\fpeval{round(
        (\zref@extractdefault{\pf@uid e}{abspage}{0}
        -\zref@extractdefault{\pf@uid a}{abspage}{0})
        +(\zref@extractdefault{\pf@uid a}{posy}{0}
        -\zref@extractdefault{\pf@uid e}{posy}{0})/\number\textheight ,2)}}%
    \protected@write\@auxout{}{\string\pf@umfangeintrag{\pf@uname}{\pf@umin}{\pf@umax}{\pf@ist}}%
    \par\smallskip\noindent
    \pf@setfarbe{\pf@umin}{\pf@umax}{\pf@ist}%
    \fcolorbox{\pf@farbe}{white}{\scriptsize
      Umfang \emph{\pf@uname}: $\approx$\,\pf@ist~Seiten
      (Soll: \pf@umin\,--\,\pf@umax~Seiten, ohne Titel/Abb./Tab./Legenden)}%
    \par\smallskip
  \fi}

%% --- Kontrollseite (nur im Entwurf; ersetzt das grüne Kapitel
%%     "Hinweise zur Arbeit mit dieser Datei", das abgegeben nicht vorkommen darf)
\newcommand{\kontrollseite}{%
  \ifAbgabe\else
  \chapter*{Kontrollseite (erscheint nur im Entwurf)}
  {\color{darkgreen}
  Modus: \textbf{Entwurf}. Für die Abgabe in \texttt{einstellungen.tex}
  \texttt{\textbackslash Abgabetrue} setzen. Dann wird alles Grüne schwarz bzw.\
  entfernt (Regel: \enquote{Abgaben mit grüner Farbe oder stehen gebliebenen
  Hilfen ergeben einen Abzug}).\par\medskip
  \textbf{Offene Platzhalter} (letzter Durchlauf): \pf@platzhalteralt\par\medskip
  \textbf{Umfangskontrolle} (Werte aus dem letzten Kompilierdurchlauf;
  Hilfstexte stehen außerhalb der Messung):\par\smallskip
  \ifx\pf@umfangliste\@empty
    \emph{noch keine Messwerte -- nochmals kompilieren.}
  \else
    {\small\begin{tabular}{p{0.55\textwidth}cc}
      \toprule Teil & Soll [Seiten] & Ist [Seiten]\\\midrule
      \pf@umfangliste
      \bottomrule
    \end{tabular}}
  \fi\par\medskip
  %% Checkliste – erscheint nur auf der Kontrollseite im Entwurf.
%% Quellen: Portfolios im Projektpraktikum 2026-2027 (Marti, 09.07.2026),
%% offizielle Gerüste 2026-2027, Zitierregeln (31.08.2026),
%% Abbildungen und Gleichungen (Limmer 2007), Moodle-Terminliste.

\textbf{Fristen (Moodle; nächste Abgabe wird erst nach der vorherigen freigeschaltet!)}\par\smallskip
{\small\begin{tabular}{@{}ll@{}}
21.\,10.\,2026 & Motivationsformular (ohne Reflexion) -- obligatorisch\\
28.\,10.\,2026 & Arbeitsplan + Reflexion -- obligatorisch\\
11.\,11.\,2026 & Wissenschaftliche Arbeit + Reflexion -- obligatorisch\\
09.\,12.\,2026 & Protokoll eines Praktikumstages + Reflexion -- obligatorisch\\
20.\,01.\,2027 & Messergebnis + Reflexion -- obligatorisch\\
23.\,01.\,2027 & Späteste Anmeldung mündliche Leistung (Portal)\\
27.\,01.\,2027 & Folien des Vortrags + Materialien für Mitstudierende\\
28.\,01.\,2027 & Vorträge (Prüfung, Anwesenheitspflicht)\\
07.\,02.\,2027 & Folien mit Kommentar + Reflexion (vorgeschlagen, freiwillig)\\
21.\,02.\,2027 & Aussage aus anderem Vortrag + Reflexion (vorgeschlagen, freiwillig)\\
10.\,03.\,2027 & Späteste Anmeldung schriftliche Leistung (Portal)\\
\textbf{15.\,03.\,2027} & \textbf{Endabgabe Portfolio -- verspätet = nicht bestanden}\\
\end{tabular}}\par\smallskip
{\small Praktikumstage (Google Kalender): donnerstags 10--12\,Uhr, O26/233 bzw.\ O26/234 --
15.\,10.\ bis 17.\,12.\,2026 und 14.\,01.\ bis 11.\,02.\,2027; keine Termine 18.\,12.--13.\,01.
Die Anmeldetermine im Hochschulportal (23.\,01./10.\,03.) stammen aus dem Plan von Petros -- im Portal prüfen!}
\par\medskip

\textbf{Vor jeder Abgabe prüfen}
\begin{itemize}\itemsep0pt
\item \texttt{\textbackslash Abgabetrue} gesetzt, keine grüne Farbe, keine Hilfen, \emph{0 Platzhalter}.
\item Deckblatt vollständig: Name, Vorname, Matrikelnummer, Studiengang, E-Mail, Thema (Projekt + Schwerpunkt).
\item Inhaltsverzeichnis mit Seitenzahlen, auch der Anhänge; durchgehend Seitenzahlen.
\item Längen eingehalten (siehe Tabelle oben; Abbildungen/Tabellen/Legenden zählen nicht mit):
      Einleitung 2--4, Beschreibung $\tfrac14$--$\tfrac12$, Diskussion der Physik 1--3,
      Reflexion 1--2, Gesamtreflexion 3--6 Seiten. Schriftgröße 12\,pt (fest eingestellt).
\item Beschreibung und Reflexion sind \emph{keine Nacherzählung}; die Leitfragen werden in Fließtext
      beantwortet (nicht in Einzelsätzen), erkennbar zugeordnet.
\item Paper (Einlage 2) wird in den späteren Einlagen zur Argumentation verwendet; Gleichungen werden
      zur Argumentation genutzt.
\item Peer-Feedback (Kommilitonen) eingeholt, aufgegriffen -- und zitiert. Petros' Arbeitsplan als Vorlage
      zitieren (Key \texttt{Petros2026Arbeitsplan}); Schwerpunkt-Abgrenzung mit Janosch festhalten.
\item Jede Übernahme zitiert: Literatur, Bilder (\texttt{\textbackslash cite[Abb.~3.4]\{key\}}),
      Gleichungen (\texttt{\textbackslash cite[Gl.~(3)]\{key\}}), Aussagen von Personen, Anleitungen,
      Datenblätter, Webseiten (mit Abrufdatum), \textbf{jede KI-Nutzung mit System und ALLEN Prompts}.
\item Journalartikel: Autoren, Titel, Journal, Band, Seiten, Jahr, DOI (nicht als Link).
\item Abbildungen: Bildunterschrift, im Text referenziert (\texttt{\textbackslash ref}), Schrift
      $\geq 10$\,pt, Achsen mit Einheit als \enquote{$f$ / kHz} (nicht \enquote{[kHz]}), ggf.\ Residuenplot.
\item Tabellen: Überschrift, im Text referenziert; Messdaten nicht als lange Listen im Text, sondern im Anhang.
\item Wichtige Gleichungen nummeriert, alle Symbole definiert und einheitlich verwendet.
\item Dateiname für Moodle: \texttt{Manevski\_WelcheEinlage.pdf} (nur PDF oder DOCX).
\item Bis auf den Arbeitsplan müssen sich alle Einlagen von denen der Partnerin/des Partners unterscheiden.
\end{itemize}

  }%
  \fi}
\providecommand{\enquote}[1]{\glqq #1\grqq}

%% --- Anhang aus PDF mit durchgehenden Seitenzahlen (Regel: "Die Arbeit ist
%%     durchgängig mit Seitenzahlen versehen ... auch der Anhänge")
%%   \anhangpdf{Kapiteltitel}{label}{datei.pdf}           alle Seiten
%%   \anhangpdf[pages={1-3}]{Kapiteltitel}{label}{datei.pdf}
%%   Die Kapitelüberschrift steht auf einer eigenen Seite, danach folgen die
%%   PDF-Seiten (auf 82 % verkleinert, damit Kopf-/Fußzeile frei bleiben).
\newcommand{\anhangpdf}[4][pages=-]{%
  \chapter{#2}\label{#3}%
  \includepdf[#1,scale=0.82,pagecommand={\thispagestyle{scrheadings}}]{#4}}
%% Gleiches ohne neues Kapitel (für weitere PDFs im selben Anhang)
\newcommand{\pdfseiten}[2][pages=-]{%
  \includepdf[#1,scale=0.82,pagecommand={\thispagestyle{scrheadings}}]{#2}}

\makeatother

%% Antwortfeld im Motivationsformular (Anhang der Endabgabe)
\newcommand{\antwort}[1]{\par\smallskip #1\par\bigskip}

\ohead{\zubearbeiten{\Projekt}}

%-----------------------------------------------------------------------------
\begin{document}

%% Titelseite der ENDABGABE – wörtlich aus pp-portfolio-geruest-2026-2027.tex
%% Hier nichts eintragen: alle Angaben kommen aus einstellungen.tex
\title{Universit\"{a}t Ulm\\\vspace{1cm}
{\LARGE \textnormal{\textsc{Fakult\"{a}t f\"{u}r Naturwissenschaften}}}\\\vspace{2cm} {\LARGE \textnormal{\textsl{Projektpraktikum \Semester}}}\\\vspace{3cm} \zubearbeiten{\textit{\Projekt}}\\
{\Large \zubearbeiten{\textit{\Schwerpunkt}}}}

\author{{\LARGE \zubearbeiten{\IhrName}}\\ $_{}$\\$_{}$\\\vspace*{1.cm}
Matrikelnummer:  \zubearbeiten{\Matrikelnummer}\\\vspace{0.8cm} Studiengang: \zubearbeiten{\Studiengang}\\\vspace{1cm} e-Mail: \zubearbeiten{\emailAdresse}}

\date{\zubearbeiten{\dasDatum}}

\maketitle

%-----------------------------------------------------------------------------
\tableofcontents
\kontrollseite   % nur im Entwurf sichtbar
% ----------------------------------------------------------------------------

%%%%%%%%%%%%%%%%%%%%%%%%%%%%%%%%%%%%%%%%%%%%%%%%%%%%%%%%%%%%%%%%%%%%%%%%%%%%%%%%
%%  EINLEITUNG  (wird ZULETZT geschrieben)   Umfang: 2–4 Seiten
%%  Gewicht 15 % der Note. Grundlage: dein Motivationsformular (Anhang).
%%%%%%%%%%%%%%%%%%%%%%%%%%%%%%%%%%%%%%%%%%%%%%%%%%%%%%%%%%%%%%%%%%%%%%%%%%%%%%%%
\chapter{Einleitung}
\label{chap:einleitung}

\begin{hinweise}
In diesem Abschnitt erörtern Sie die folgenden Themen (und nicht nur in einem Satz!):
\begin{itemize}\itemsep0pt
\item Motivation für die Wahl Ihres Projektthemas.
\item Motivation für die Wahl Ihres Schwerpunktthemas innerhalb Ihres Projektthemas.
\item Wie unterscheidet sich Ihr Projektthema von anderen Themen im aktuellen Semester?
\item Wie unterscheidet sich Ihr Projektthema von ähnlichen Themen in früheren Jahren?
\item Wie unterscheidet sich Ihr Schwerpunktthema von dem des anderen Gruppenmitglieds? (wenn Sie eine Zweiergruppe bilden)
\item Warum wollen Sie das von Ihnen gewählte Schwerpunktthema bearbeiten?
\item Was ist der Kontext Ihres Schwerpunktthemas in Ihrem Projekt?
\item Welche Anregungen hatten Sie?
\item Was erwarten Sie von der Bearbeitung Ihres Schwerpunktthemas?
\item Wie steht es mit Ihren Vorkenntnissen in Bezug zum Projektthema?
\item Wie steht es mit Ihren Vorkenntnissen in Bezug zum Schwerpunktthema?
\item Gibt es weitere Punkte, die Sie zu Ihren Vorkenntnissen anmerken wollen?
\item Weitere für Sie wichtige Punkte zu diesem Portfolio.
\item Vorschau über das Portfolio bestehend aus einer Auflistung der einzelnen Bestandteile sowie einer kurze Einführung in deren Inhalte.
\end{itemize}
Sie können dazu das Motivationsformular (als Anlage beilegen) als Grundlage verwenden.
\textbf{Die Einleitung ist mindestens zwei Seiten und maximal vier Seiten lang. Die Längenangaben sind ohne Titel, Bilder, Tabellen und Tabellen- oder Bildlegenden zu verstehen.}

\smallskip\emph{Zusätzlich aus dem Leitfaden \enquote{Portfolios im Projektpraktikum} und den Bewertungskriterien:}
Welche Gründe aus Ihrer Lernbiografie haben Sie zu dem Thema geführt? Zusammenhang Schwerpunktthema
\(\leftrightarrow\) Projektthema in mehr als einem Satz. Schwerpunkt von dem der Partnerin/des Partners abgrenzen:
\emph{Was außer dem Titel ist der Unterschied?} Das Schwerpunktthema muss klar definiert sein (fehlt es, gibt es Abzug).
Tipp: Unterabschnitte (\texttt{\textbackslash section*}) helfen, die geforderten Aussagen identifizierbar zu machen.
\end{hinweise}

\begin{umfang}{Einleitung}{2}{4}

\section*{Motivation und Themenwahl}
\platzhalter{Motivation für das Projektthema; Anregungen; Lernbiografie; Vorkenntnisse (Grundlage:
Motivationsformular, Anhang~\ref{anh:motivationsformular}) -- Stichpunkte: \enquote{\Projekt} ist eine
eigene Idee der Studierenden \cite{Themenliste2026}: Ein Rattleback (Keltischer Wackelstein) ist ein
länglicher Körper mit gewölbter Unterseite, der sich in eine Richtung stabil dreht und, in die andere
Richtung angestoßen, zu nicken beginnt und seine Drehrichtung umkehrt; Ziel: diese Umkehrdynamik mit
3D-gedruckten Körpern und eingebautem Inertialsensor quantitativ untersuchen und mit der Theorie
vergleichen; scheinbarer Widerspruch zur Drehimpulserhaltung als Anreiz}

\section*{Schwerpunktthema und Abgrenzung}
\platzhalter{Definition des Schwerpunkts \Schwerpunkt; Kontext im Projekt; Abgrenzung zu
\Partner{} (\PartnerSchwerpunkt); Abgrenzung zu anderen/früheren Themen; Erwartungen -- Stichpunkte:
eigener Teil laut Arbeitsplan = Bewegungsgleichungen, numerische Simulation der Umkehrdynamik (Moden,
Energiebilanz), Inertialsensor (Firmware, Kalibrierung) und Python-Datenpipeline; Janosch: Körperparameter
und Trägheitstensor aus dem CAD-Modell, Vorhersage der Umkehrzeit in Abhängigkeit von Anfangs-$\omega$ und
$\delta$, Konstruktion und 3D-Druck; Unterschied außer dem Titel: Modell/Simulation vs.\ Körper/Vorhersage,
beide treffen sich beim Vergleich Messung $\leftrightarrow$ Theorie; andere Themen dieses Semesters
\cite{Themenliste2026} (z.\,B.\ Stabilität rollender Fahrräder, Doppelpendel, rotierender
Tischtennisball, schräge Leiter) behandeln andere Systeme -- keines eine Kopplung von Schwingung und
Drehung mit Drehrichtungsumkehr; in den Themenlisten kein Rattleback-Thema aus früheren Jahren (prüfen)}

\section*{Vorschau auf das Portfolio}
\platzhalter{Auflistung aller Bestandteile mit kurzer Einführung in die Inhalte: Arbeitsplan
(Abschnitt~\ref{sec:arbeitsplan}), Paper von Kondo und Nakanishi (Abschnitt~\ref{sec:paper}),
Laborjournal-Tag (Abschnitt~\ref{sec:protokoll}), Messergebnis (Abschnitt~\ref{sec:messergebnis}),
eigener Vortrag (Abschnitt~\ref{sec:vortrag}), Aussage aus einem anderen Vortrag
(Abschnitt~\ref{sec:fremdvortrag}), Gesamtreflexion (Kapitel~\ref{chap:gesamtreflexion})}

\end{umfang}


%-----------------------------------------------------------------------------
\chapter{Hauptteil}
\hinweis{Der Hauptteil besteht aus mindestens sechs Einlagen. Jede Einlage: Beschreibung + Reflexion
(+ Diskussion der Physik bei Einlagen 2, 4, 6); die Einlage selbst steht im Anhang.}

%%%%%%%%%%%%%%%%%%%%%%%%%%%%%%%%%%%%%%%%%%%%%%%%%%%%%%%%%%%%%%%%%%%%%%%%%%%%%%%%
%%  EINLAGE 1: ARBEITSPLAN          Zwischenabgabe: 28.10.2026 (obligatorisch)
%%  Struktur: Beschreibung (1/4–1/2 S.) + Reflexion (1–2 S.)
%%  Die Einlage selbst (der Plan) steht im Anhang: 1_arbeitsplan_anhang.tex
%%%%%%%%%%%%%%%%%%%%%%%%%%%%%%%%%%%%%%%%%%%%%%%%%%%%%%%%%%%%%%%%%%%%%%%%%%%%%%%%
\section{Einlage: Arbeitsplan}
\label{sec:arbeitsplan}

\begin{hinweise}
Diese Einlage ist Ihr Arbeitsplan. Der Arbeitsplan ist die \emph{einzige} Einlage, die mit der Partnerin/dem
Partner identisch sein darf (\enquote{Ihr gemeinsamer Arbeitsplan}). Hinweise zur Erstellung (Marti): nicht nur
eine grobe Liste, sondern feine Teilaufgaben mit zeitlichem Ablauf, sichtbar machen, was parallel läuft;
Vorbereitung (\enquote{Einkaufen}) und Nachbereitung (\enquote{Abwasch}) nicht vergessen -- z.\,B. als
Excel-Tabelle oder in ProjectLibre, oder mit dem Gantt-Gerüst im Anhang.

\smallskip\emph{Zusätzlich aus der Moodle-Aufgabe (Frist 28.\,10.\,2026, 23:59\,Uhr):}
\begin{itemize}\itemsep0pt
\item Den Plan mit den Betreuern besprechen -- deren Rückmeldung in der Reflexion aufgreifen und als
      Gespräch zitieren (Vorlage \texttt{VORLAGE-Gespraech} in \texttt{literatur.bib}).
\item Gleicher Arbeitsplan für beide, aber die Reflexion schreibt jeder \emph{einzeln und unabhängig} von
      \Partner.
\item Reflexion ist keine Nacherzählung; Beschreibung $1/4$--$1/2$ Seite, Reflexion 1--2 Seiten
      (ohne Titel, Abbildungen, Tabellen und Legenden).
\item Arbeitsplan und Reflexion als \emph{ein PDF} hochladen (\texttt{Manevski\_Arbeitsplan.pdf}) --
      danach den \textbf{Abgabeknopf} drücken!
\end{itemize}
\end{hinweise}

\subsection{Beschreibung}
\begin{hinweise}
Die Beschreibung gibt alle notwendigen Informationen an, die die Einlage beschreiben, also zum Beispiel den
Fundort, die Literaturstelle und eine kurze Inhaltsübersicht, aber keine Nacherzählung ($1/4$ Seite, max.
$1/2$ Seite!). Leitfaden: Um welche Art von Einlage handelt es sich? Woher stammt sie? Welches sind die
wesentlichen Aussagen?
\end{hinweise}
\begin{hinweise}
\emph{Deine Gliederung aus dem Fragebogen (ca.\ $1/3$ Seite, 5--6 Sätze):}
\begin{enumerate}\itemsep0pt
\item \emph{Was:} gemeinsamer Arbeitsplan mit \Partner, in Excel als Gantt-Diagramm, Wochenraster
      KW\,42/2026 bis KW\,11/2027, im Anhang~\ref{anh:arbeitsplan}.
\item \emph{Woher:} am 22.\,10.\,2026 fertiggestellt und am selben Tag mit Prof.\ Marti besprochen
      \cite{Marti2026Gespraech}; Baumstruktur je parallel laufendem Thema nach dem Beispiel
      \enquote{Mittagessen kochen} \cite{MoodleArbeitsplan2026}; alles gemeinsam erstellt.
\item \emph{Aufbau:} sieben Blöcke (Theorie \& Literatur, Materialien, Versuchsaufbau, Messungen,
      Auswertung, Vortrag, Portfolio); jede Aufgabe mit Person, Start, Ende, Status; Prozentangabe je
      Block; alle Abgaben eingetragen.
\item \emph{Kernaussage (1.7):} man sieht sofort, welche Aufgaben parallel laufen und welche
      zusammenhängen.
\item \emph{Nutzung (1.10):} Plan bleibt in der abgegebenen Form, die Statusspalte dient zum Verfolgen
      des Fortschritts. \textbf{Achtung:} Janosch schreibt, der Plan sei eine \enquote{Arbeitshypothese,
      die wir anpassen} -- mit ihm abstimmen, damit sich eure Abgaben nicht widersprechen.
\end{enumerate}
\end{hinweise}
\begin{umfang}{Einlage 1 (Arbeitsplan): Beschreibung}{0.25}{0.5}
\noindent\emph{\enquote{Dieser Text ist von KI formuliert worden.}}~\cite{Manevski2026KIArbeitsplan}

Bei dieser Einlage handelt es sich um den Arbeitsplan für unser Projekt \enquote{\Projekt}
(Anhang~\ref{anh:arbeitsplan}). \Partner{} und ich haben ihn gemeinsam in Excel als Gantt-Diagramm
aufgebaut, im Wochenraster von KW~42/2026 bis KW~11/2027. Die Baumstruktur, in der jedes parallel laufende
Thema seine eigenen Zweige hat, haben wir uns beim Beispiel \enquote{Mittagessen kochen} aus der
Aufgabenstellung abgeschaut \cite{MoodleArbeitsplan2026}. Fertig war der Plan am 22.\,10.\,2026, am
selben Tag haben wir ihn mit Prof.~Marti besprochen \cite{Marti2026Gespraech}.

Gegliedert ist er in sieben Blöcke, von Theorie und Literatur über Materialien, Versuchsaufbau, Messungen
und Auswertung bis zu Vortrag und Portfolio. Zu jeder Aufgabe steht drin, wer sie übernimmt, wann sie
beginnt und endet und wie weit sie ist; pro Block gibt es eine Prozentanzeige, und alle Abgaben sind
markiert. Die wichtigste Aussage ist für mich, dass man auf einen Blick sieht, welche Aufgaben parallel
laufen und welche zusammenhängen. Als Meilensteine sehe ich den funktionierenden Prototyp, die fertige
Simulation und eine saubere Messreihe.
\end{umfang}

\subsection{Reflexion}
\begin{hinweise}
In der Reflexion sollen Sie sich Gedanken machen über die Wirkung dieser Einlage und der Arbeit mit ihr auf
Ihren Lernfortschritt.
\begin{itemize}\itemsep0pt
\item In welchem Zusammenhang steht der Arbeitsplan zu Ihrem Projektthema?
\item In welchem Zusammenhang steht der Arbeitsplan zu Ihrem Schwerpunktthema?
\item Was haben Sie bei der Erstellung des Arbeitsplanes gelernt?
\item Wie hat Ihnen der Arbeitsplan beim Verständnis Ihres Projektthemas und Ihres Schwerpunktthemas geholfen?
\item Wo gab es Schwierigkeiten bei der Beschäftigung mit der Arbeitsplan?
\item Was sagen Ihre Kommilitoninnen oder Kommilitonen dazu, als Kommentar oder wenn Sie mit Ihnen über den Arbeitsplan diskutieren?
\item Enthält Ihr Arbeitsplan logische oder physikalische Fehler?
\item Gab es bei der Beschäftigung mit dem Arbeitsplan eine überraschende neue Erkenntnis?
\end{itemize}
\textbf{Die Reflexion für eine Einlage sollte mindestens eine Seite und maximal 2 Seiten lang sein. Die
Längenangaben sind ohne Titel, Bilder, Tabellen und Tabellen- oder Bildlegenden zu verstehen.}
Peer-Feedback zitieren (siehe \texttt{literatur.bib}, Vorlage \texttt{VORLAGE-Gespraech}).

\smallskip\emph{Deine Gliederung aus dem Fragebogen (Ziel ca.\ 1{,}2 Seiten, vier Absätze, Ich-Form,
eher sachlich mit persönlichen Stellen):}
\begin{enumerate}\itemsep0pt
\item \emph{Wie der Plan entstanden ist (A1--A4, D1, D3, D4):} anfangs nicht in Worte zu fassen, was
      wir wollten; Brainstorming und gezeichnete \enquote{Prototypen}; mehrere Runden im Kopf und auf
      Papier, erst dann grob in Excel. Zu jedem Schwerpunkt alle Aufgaben gesammelt, aufgeteilt, nach
      Zusammenhang und sinnvollem Zeitabschnitt gruppiert, Dauer diskutiert. Geholfen hat die Erfahrung
      aus zwei Jahren Grundpraktikum: \enquote{schnelle} Aufgaben sind nie schnell erledigt. Am Anfang
      Meinungsverschiedenheiten mit \Partner, gelöst über Argumente und vergessene Themen. Gelernt: wie
      wichtig Iterieren ist -- mit Abstand sieht man, was fehlt, und hinterfragt Erarbeitetes. Noch nie
      einen so detaillierten Plan gemacht.
\item \emph{Projektthema (B1--B4, G7):} vorher nur grob bekannt, wie der Stein funktioniert. Erst beim
      Ableiten der Aufgaben aus dem Versuchsziel wurde klar, was in der Theorie steckt. Erst beim Planen
      als eigene Schritte erkannt: Reibung, reproduzierbarer Anstoß, Sensorkalibrierung. Aufwendiger als
      gedacht: Theorie und Vorbereitung, das 3D-Modell muss für den kleinen Effekt sehr genau sein, und
      schon kleinste Änderungen der Startbedingungen spielen eine große Rolle. Offen: ob Glasplatte und
      Anstoß das überhaupt reproduzierbar erlauben.
\item \emph{Schwerpunkt (C2--C5, H2):} eigene Zeilen: Bewegungsgleichungen (1.4, 29.\,10.--25.\,11.),
      Firmware und Kalibrierung (2.1.3/2.1.4, 05.--26.\,11.), Simulation (1.6, 19.\,11.--17.\,12.),
      Datenpipeline (5.1, 26.\,11.--17.\,12.), Video (3.2). \textbf{Aus dem Plan abgelesen:} vom 19.\ bis
      25.\,11.\ laufen 1.4, 1.6, 2.1.3 und 2.1.4 gleichzeitig -- dein Engpass; die Simulation startet
      am 19.\,11., die Körperparameter von Janosch (1.5) sind erst am 25.\,11.\ fertig -- deine direkte
      Abhängigkeit von ihm (alle wichtigen Größen des Körpers). Motivation: Programmierung, Theorie und
      Simulation zwingen dazu, sich um kleinste Details zu kümmern, die sonst im Verständnis fehlen
      würden. Erst wenig Lust aufs Projekt erwartet, nach dem Ausarbeiten der eigenen Aufgaben jetzt Spaß.
\item \emph{Schwierigkeiten, Rückmeldungen, Ausblick (D2, E1--E3, F2--F4, G, I1, I4):} der Gedanke an
      so einen detaillierten Plan war anfangs einschüchternd, danach eher Zufriedenheit. Neigung, Zeit
      zu unterschätzen -- bekannt und bewusst mit Puffer kompensiert, v.\,a.\ beim 3D-Druck (Teillieferungen)
      und weil wir vom Betreuer abhängen, der den Stein druckt. Prof.\ Marti hat Fragen gestellt und den Plan
      dann so akzeptiert \cite{Marti2026Gespraech}; von Kommilitonen nur, dass er anfangs hässlich war --
      optisch überarbeitet, Struktur gleich gelassen. Auffällige Zeitüberschneidungen geprüft: gewollt als
      Puffer bzw.\ weil die Messungen selbst kurz sind. Kontrolle über ständige Absprache und die
      Statusspalte. Was ich beim nächsten Plan anders mache, kann ich erst beurteilen, wenn ich versucht
      habe, diesen einzuhalten.
\end{enumerate}
\emph{Unabhängig von Janosch bleiben} (Moodle: Reflexionen \enquote{einzeln und unabhängig}): nicht seine
Absatzthemen übernehmen (Messtechnik/Puffer, Messkonzept Messreihe~3, Detailgrad) und nicht seine
Formulierungen (\enquote{vom Ende her denken}, \enquote{Arbeitshypothese}).
\end{hinweise}
\begin{umfang}{Einlage 1 (Arbeitsplan): Reflexion}{1}{2}
\noindent\emph{\enquote{Dieser Text ist von KI formuliert worden.}}~\cite{Manevski2026KIArbeitsplan}

Am Anfang konnten Janosch und ich gar nicht richtig in Worte fassen, wie der Plan aussehen soll. Wir haben
deshalb erst gebrainstormt und ein paar \enquote{Prototypen} auf Papier gezeichnet; der Plan ist also
zuerst im Kopf entstanden, dann in mehreren Runden auf dem Blatt und erst zum Schluss in Excel. Da unsere
Schwerpunkte schon feststanden, haben wir zu jedem Schwerpunkt alle Aufgaben aufgeschrieben, aufgeteilt und
danach sortiert, was zusammenhängt und in welchen Abschnitt des Praktikums es passt. Am längsten haben wir
darüber diskutiert, wie viel Zeit jede Aufgabe bekommt. Geholfen hat mir dabei die Erfahrung aus dem
Grundpraktikum: Aufgaben, bei denen man denkt, dass sie schnell erledigt sind, sind eigentlich nie schnell
erledigt. Anfangs waren
wir uns nicht überall einig, aber mit Argumenten und Themen, die einer von uns vergessen hatte, sind wir bei
einem Plan gelandet, mit dem wir beide zufrieden sind. So einen detaillierten Plan hatte ich vorher noch nie
gemacht. Gelernt habe ich vor allem, wie viel das Iterieren bringt: Schaut man nach einer Weile mit anderem
Blick auf das schon Erarbeitete, merkt man, was fehlt, und fragt sich, ob alles wirklich stimmt und wichtig ist.

Über den Wackelstein wusste ich vorher eigentlich nur, dass es ihn gibt und wie er grob funktioniert. Erst
als wir aus dem Ziel des Versuchs die einzelnen Aufgaben abgeleitet haben, ist mir klar geworden, wie viel
Theorie hinter dem Effekt steckt, die man beim bloßen Anschauen des Steins nicht sieht. Die Reibung auf der
Glasplatte, einen reproduzierbaren Anstoß oder die Kalibrierung des Sensors hatte ich nie als eigene
Arbeitsschritte gesehen, dabei sind genau das die Stellen, an denen der Versuch scheitern kann. Aufwendiger
als gedacht sind vor allem Theorie und Vorbereitung: Für so einen kleinen Effekt muss das 3D-Modell praktisch
perfekt sein, und schon kleinste Änderungen der Startbedingungen spielen eine große Rolle. Ob die Glasplatte
Rollen ohne Gleiten zulässt und ob sich der Anstoß reproduzieren lässt, müssen wir noch herausfinden.

Mein Schwerpunkt sind die Bewegungsgleichungen und die Simulation der Spin-Umkehr. Im Plan gehören dazu die
Bewegungsgleichungen (1.4), Firmware und Kalibrierung des Sensors (2.1.3, 2.1.4), die Simulation (1.6) und
die Datenpipeline in Python (5.1). Beim Einteilen dieser Aufgaben habe ich besser verstanden, was mein
Schwerpunkt konkret bedeutet und womit er zusammenhängt. Vom 19.\ bis 25.\,11.\ laufen zum Beispiel
Bewegungsgleichungen, Simulation, Firmware und Kalibrierung gleichzeitig, das wird meine vollste Woche.
Außerdem brauche ich für die Simulation alle wichtigen Größen des Körpers von Janosch (1.5), die laut Plan
erst am 25.\,11.\ fertig sind, während die Simulation schon am 19.\,11.\ beginnt. An dieser Stelle hängt mein
Teil am stärksten von seinem ab. Gereizt haben mich an diesen Aufgaben das Programmieren, die Theorie und die
Simulation, weil man sich dabei um jedes Detail kümmern muss, an das man ohne dieses Ziel nie denken würde
und das dann im Verständnis fehlt. Ehrlich gesagt hatte ich erwartet, dass mir das Projekt wenig Spaß macht;
seit ich meine Aufgaben ausgearbeitet habe, sehe ich das anders.

Wirklich schwer gefallen ist mir beim Planen nichts, auch wenn mich der Gedanke an so einen detaillierten
Plan am Anfang eher eingeschüchtert hat. Ich weiß von mir, dass ich die Zeit für Aufgaben gern unterschätze.
Das haben wir bewusst mit viel Puffer ausgeglichen, besonders beim 3D-Druck und bei Teillieferungen, wo wir
zusätzlich von unserem Betreuer abhängen, der den Stein druckt. Prof.~Marti hat zu dem Plan einige Fragen
gestellt, die wir beantworten konnten, und ihn danach so akzeptiert \cite{Marti2026Gespraech}. Von
Kommilitonen kam vor allem, dass der Plan anfangs nicht schön aussah; wir haben ihn daraufhin optisch
überarbeitet, an der Struktur aber nichts geändert. Logische Fehler sehe ich keine: Wo sich Zeiträume
überschneiden, ist das gewollt, als Puffer für Nachmessungen oder weil die Messungen selbst nicht lange
dauern. Auf dem Laufenden halten wir uns über ständige Absprache und die Statusspalte. Was ich beim nächsten
Plan anders machen würde, kann ich erst sagen, wenn ich versucht habe, diesen hier einzuhalten.
\end{umfang}

%% Optional (vom Leitfaden ausdrücklich erwünscht): eigene Bewertungskriterien
%\subsection{Eigene Bewertungskriterien}
%Für diese Einlage schlage ich zusätzlich folgende Kriterien vor: ...

%%%%%%%%%%%%%%%%%%%%%%%%%%%%%%%%%%%%%%%%%%%%%%%%%%%%%%%%%%%%%%%%%%%%%%%%%%%%%%%%
%%  EINLAGE 2: WISSENSCHAFTLICHE ARBEIT     Zwischenabgabe: 11.11.2026 (obligatorisch)
%%  Struktur: Beschreibung (1/4–1/2 S.) + Diskussion der Physik (1–3 S.) + Reflexion (1–2 S.)
%%  Paper: Kondo & Nakanishi, Phys. Rev. E 95, 062207 (2017) -> \PaperKey = KondoNakanishi2017.
%%  Janosch darf NICHT dasselbe Paper nehmen. Falls doch: auf Garcia & Hubbard (1988),
%%  Key GarciaHubbard1988, wechseln (liegt ebenfalls in Literatur/).
%%%%%%%%%%%%%%%%%%%%%%%%%%%%%%%%%%%%%%%%%%%%%%%%%%%%%%%%%%%%%%%%%%%%%%%%%%%%%%%%
% \PaperKey wird in einstellungen.tex gesetzt

\section{Einlage: Eine wissenschaftliche Arbeit zu einem Teilaspekt Ihres Schwerpunktgebietes}
\label{sec:paper}

\begin{hinweise}
Diese Einlage muss aus einer begutachteten Fachzeitschrift sein. Sie sollen mit dieser Einlage die
professionelle Literaturrecherche lernen. Bücher, Konferenzbeiträge, Trade Journals und ähnliches werden hier
nicht akzeptiert.
Dieser Artikel \textbf{muss} in Ihrer weiteren Arbeit verwendet werden (v.\,a. in Einlage 4: \enquote{Eine
Bezug zu Ihrem Paper ist Voraussetzung}).
Sie dürfen nicht den gleichen Artikel verwenden wie Ihre Partnerin / Ihr Partner.
Akzeptable wissenschaftliche Einlagen sind begutachtete wissenschaftliche Artikel aus wissenschaftlichen
Zeitschriften, die Sie über webofknowledge.com oder auch über Google Scholar finden. Erfüllt die Einlage die
Anforderung nicht, wird sie in der Bewertung nicht berücksichtigt. Bei Zweifeln über die Zulässigkeit der
beiden Einlagen mit wissenschaftlichen Arbeiten fragen Sie Prof.~O.~Marti.
\end{hinweise}

\subsection{Beschreibung}
\begin{hinweise}
Die Beschreibung gibt alle notwendigen Informationen an, die die Einlage beschreiben, also zum Beispiel den
Fundort, die Literaturstelle und eine kurze Inhaltsübersicht, aber keine Nacherzählung ($1/4$ Seite, max.
$1/2$ Seite!). Leitfaden: Art der Einlage, wie bin ich darauf aufmerksam geworden (Datenbank, Suchbegriffe,
Hinweis von \dots), wesentliche Aussagen in eigenen Worten.
\end{hinweise}
\begin{umfang}{Einlage 2 (Paper): Beschreibung}{0.25}{0.5}
\platzhalter{In eigenen Worten ausformulieren -- Stichpunkte: Y.~Kondo und H.~Nakanishi,
\enquote{Rattleback dynamics and its reversal time of rotation}, Phys.\ Rev.\ E \textbf{95}, 062207
(2017) \cite{\PaperKey}, begutachtete Fachzeitschrift (APS); Fundort: Literaturrecherche in Web of
Science/Google Scholar (Arbeitspaket 1.1), Suchbegriffe: \dots; Inhalt: physikalisch transparente
Neuformulierung der Rattleback-Dynamik, Reduktion auf drei Variablen (Drehen, Nicken, Rollen),
Garcia-Hubbard-Formel für die Umkehrzeit als Produkt von vier Faktoren, ausführliche numerische
Simulationen zu Gültigkeit und Grenzen der Formel; PDF in Anhang~\ref{anh:paper}}
\end{umfang}

\subsection{Diskussion der Physik}
\begin{hinweise}
Beschreiben Sie hier die Physik dieser Einlage. Fassen Sie den Artikel zusammen. Arbeiten Sie besonders die
Physik, die für ihr Schwerpunktthema relavant ist, heraus.
\begin{itemize}\itemsep0pt
\item Was ist die Hauptaussage der wissenschaftlichen Arbeit?
\item Was sind die Voraussetzungen für den Inhalt des Artikels?
\item Wie hilft physikalisch die Physik dieser Einlage bei Ihrem Schwerpunktthema?
\item Wo gab es Schwierigkeiten beim Verständnis der Physik? Welche? Mit welcher zusätzlichen Literatur haben Sie das Problem gelöst
\item Haben Sie weitere Artikel oder Bücher konsultiert um das Paper zu verstehen? Welche?
\item Welche Aussagen verwenden Sie für Ihr Schwerpunktthema und wo? Warum?
\end{itemize}
\textbf{Die Diskussion der Physik sollte mindestens eine Seite und maximal 3 Seiten lang sein. Die
Längenangaben sind ohne Titel, Bilder, Tabellen und Tabellen- oder Bildlegenden und ohne übergrosse Abstände
zu Gleichungen zu verstehen.}
Gleichungen zur Argumentation verwenden, nummerieren, Symbole definieren; aus dem Paper übernommene Gleichungen
mit Gleichungsnummer zitieren: \texttt{\textbackslash cite[Gl.~(3)]\{\PaperKey\}}.
\end{hinweise}
\begin{umfang}{Einlage 2 (Paper): Diskussion der Physik}{1}{3}
\platzhalter{Hauptaussage, Voraussetzungen, relevante Physik für \Schwerpunkt{} -- Stichpunkte aus dem
Abstract/der Einleitung des Papers:
(1) Drei Voraussetzungen für die Spin-Umkehr: unterschiedliche Hauptkrümmungen der Unterseite,
unterschiedliche horizontale Hauptträgheitsmomente, Fehlstellung (Winkel $\delta$) zwischen
Trägheits- und Krümmungshauptachsen;
(2) die Fehlstellung koppelt Drehen mit Nicken und Rollen -- asymmetrisch, mit entgegengesetztem
Vorzeichen des Drehmoments, daher Umkehr je nach Drehsinn unterschiedlich schnell;
(3) Annahmen: Rollen ohne Gleiten, keine Dissipation;
(4) Umkehrzeit $t_r$ (Garcia-Hubbard \cite{GarciaHubbard1988}) als Produkt aus Fehlstellungswinkel,
Differenz der inversen Trägheitsmomente, Differenz der Krümmungsradien und quadrierter
Schwingungsfrequenz -- Gleichung mit Nummer aus dem Paper zitieren;
(5) Simulationen: Formel gut bei kleiner Drehung und kleinen Schwingungen, bei schneller Drehung
keine Umkehr, Taumeln oder chaotisches Verhalten;
Bezug zum eigenen Schwerpunkt: Bewegungsgleichungen (Arbeitspaket 1.4) und eigene Simulation (1.6)
gegen die Ergebnisse des Papers prüfen}

% Beispiel für eine zitierte, nummerierte Gleichung (anpassen oder löschen):
% Nach \cite[Gl.~(1)]{\PaperKey} gilt
% \begin{equation}
%   t_r = \dots     % Garcia-Hubbard-Formel in der Form von Kondo & Nakanishi
%   \label{eq:paper-hauptgleichung}
% \end{equation}
% wobei ... ist. Aus Gl.~\eqref{eq:paper-hauptgleichung} folgt ...

\end{umfang}

\subsection{Reflexion}
\begin{hinweise}
In der Reflexion sollen Sie sich Gedanken machen über die Wirkung dieser Einlage und der Arbeit mit ihr auf
Ihren Lernfortschritt.
\begin{itemize}\itemsep0pt
\item In welchem Zusammenhang steht die Einlage mit Ihrem Schwerpunktthema?
\item Was haben Sie gelernt?
\item Wie hat Ihnen die Einlage beim Verständnis Ihres Projektes geholfen?
\item Wo gab es Schwierigkeiten bei der Beschäftigung mit der Einlage?
\item Hätte ein schon bestandenes Modul beim Verständnis helfen können? Welches? Warum oder warum nicht?
\item Was haben Sie durch diese Zusatzliteratur gelernt?
\item Was hätten Sie anders gemacht, wenn Sie die Einlage erstellt hätten?
\item Was sagen Ihre Kommilitoninnen oder Kommilitonen dazu, als Kommentar oder wenn Sie mit Ihnen über die Einlage diskutieren?
\item Haben Sie in dem Artikel der Einlage logische oder physikalische Fehler entdeckt?
\item Gab es bei der Beschäftigung mit der Einlage eine überraschende neue Erkenntnis?
\end{itemize}
\textbf{Die Reflexion für eine Einlage sollte mindestens eine Seite und maximal 2 Seiten lang sein. Die
Längenangaben sind ohne Titel, Bilder, Tabellen und Tabellen- oder Bildlegenden zu verstehen.}
\end{hinweise}
\begin{umfang}{Einlage 2 (Paper): Reflexion}{1}{2}
\platzhalter{Reflexion in Fließtext -- Leitfragen erkennbar beantworten}
\end{umfang}

%\subsection{Eigene Bewertungskriterien}
%...

%%%%%%%%%%%%%%%%%%%%%%%%%%%%%%%%%%%%%%%%%%%%%%%%%%%%%%%%%%%%%%%%%%%%%%%%%%%%%%%%
%%  EINLAGE 3: PROTOKOLL EINES PRAKTIKUMSTAGES   Zwischenabgabe: 09.12.2026 (obligatorisch)
%%  Struktur: Beschreibung (1/4–1/2 S.) + Reflexion (1–2 S.)
%%  Die Laborjournal-Seiten (unredigiert, ALLE des Tages) kommen in den Anhang.
%%%%%%%%%%%%%%%%%%%%%%%%%%%%%%%%%%%%%%%%%%%%%%%%%%%%%%%%%%%%%%%%%%%%%%%%%%%%%%%%
% \Protokolldatum und \Laborbuchlink werden in einstellungen.tex gesetzt

\section{Einlage: Alle Originalseiten ihres Protokolls eines Praktikumstages.}
\label{sec:protokoll}

\begin{hinweise}
Diese Einlage muss aus Ihrem elektronischen Laborjournal sein. Sie verwenden alles was Sie an dem gewählten Tag
aufgeschrieben haben. Sie geben den Link zu Ihrem Laborbucheintrag an.
Verwenden Sie nicht den gleichen Tag wie Ihre Partnerin / Ihr Partner. Machen Sie sich Gedanken über die
folgenden Fragen: Was ist ein gutes Protokoll? Ein gutes Protokoll enthält alle Informationen (Gleichungen,
Literaturhinweise, Handlungsanweisungen, Schaltpläne, Konstruktionszeichnungen, usw.) die notwendig ist, dass
eine andere Person mit Ihrem Bildungsstand, aber ohne weitere Verbindungen zu Ihnen, Ihre Experimente oder
Rechnungen wiederholen kann und Ihre Ergebnisse bewerten kann.
Diese Einlage ist keine Nacherzählung Ihrer Arbeiten an einem Praktikunstag, sondern eine Analyse Ihrer
Dokumentation Ihrer Arbeiten.

\smallskip\emph{Projektbezug:} Praktikumstage sind donnerstags 10--12\,Uhr (O26/233 bzw.\ O26/234).
Abgabe am 09.\,12.\,2026, daher kommen die Tage bis 03.\,12.\ in Frage, z.\,B.\ Platine abholen und
Funktionstest (05.\,11.), Firmware/Kalibrierung des Inertialsensors (12.--26.\,11.) oder die Testmessung
des kompletten Aufbaus (03.\,12.). Tag mit Janosch absprechen (muss verschieden sein), dann
\texttt{\textbackslash Protokolldatum} und \texttt{\textbackslash Laborbuchlink} in \texttt{einstellungen.tex}
setzen.
\end{hinweise}

\subsection{Beschreibung}
\begin{hinweise}
Die Beschreibung gibt alle notwendigen Informationen an, die die Einlage beschreiben, also zum Beispiel das
Datum, die Art des Laborjournals und eine kurze Inhaltsübersicht, aber keine Nacherzählung ($1/4$ Seite, max.
$1/2$ Seite!).
\end{hinweise}
\begin{umfang}{Einlage 3 (Protokoll): Beschreibung}{0.25}{0.5}
Die Einlage umfasst alle unredigierten Seiten meines elektronischen Laborjournals vom \Protokolldatum{}
(Anhang~\ref{anh:protokoll}). Link zum Laborbucheintrag: \expandafter\url\expandafter{\Laborbuchlink}.

\platzhalter{Art/Software des elektronischen Laborjournals (gemeinsam mit \Partner, Arbeitspaket 4.1),
kurze Inhaltsübersicht des Tages (keine Nacherzählung)}
\end{umfang}

\subsection{Reflexion}
\begin{hinweise}
In der Reflexion sollen Sie sich Gedanken machen über die Wirkung dieser Einlage und der Arbeit mit ihr auf
Ihren Lernfortschritt.
\begin{itemize}\itemsep0pt
\item In welchem Zusammenhang steht die Einlage mit Ihrem Schwerpunktthema und Ihrem Projektthema?
\item Verstehen Sie noch Ihren Notizen?
\item Was haben Sie bei der Erstellung dieser Einlage gelernt?
\item Wie hat Ihnen die Einlage beim Verständnis Ihres Projektthemas geholfen?
\item Wo gab es Schwierigkeiten bei der Beschäftigung mit der Einlage?
\item Was wünschen sie sich von Ihren Notizen beim erneuten Lesen?
\item Was hätten Sie anders gemacht, wenn Sie die Einlage nochmals erstellen müssten?
\item Was sagen Ihre Kommilitoninnen oder Kommilitonen dazu, als Kommentar oder wenn Sie mit Ihnen über die Einlage diskutieren?
\item Haben Sie in Ihrem Protokoll logische oder physikalische Fehler entdeckt?
\item Gab es bei der Beschäftigung mit der Einlage eine überraschende neue Erkenntnis?
\end{itemize}
\textbf{Die Reflexion für eine Einlage sollte mindestens eine Seite und maximal 2 Seiten lang sein. Die
Längenangaben sind ohne Titel, Bilder, Tabellen und Tabellen- oder Bildlegenden zu verstehen.}
Tipp: Konkrete Stellen im Journal benennen (\enquote{Seite 2 oben fehlt die Verstärkereinstellung \dots}).
\end{hinweise}
\begin{umfang}{Einlage 3 (Protokoll): Reflexion}{1}{2}
\platzhalter{Analyse der eigenen Dokumentation: Was ist ein gutes Protokoll? Wäre mein Experiment
wiederholbar? Leitfragen erkennbar beantworten}
\end{umfang}

%\subsection{Eigene Bewertungskriterien}
%...

%%%%%%%%%%%%%%%%%%%%%%%%%%%%%%%%%%%%%%%%%%%%%%%%%%%%%%%%%%%%%%%%%%%%%%%%%%%%%%%%
%%  EINLAGE 4: EINES IHRER ERGEBNISSE     Zwischenabgabe: 20.01.2027 (obligatorisch)
%%  Struktur: Beschreibung (1/4–1/2 S.) + Diskussion der Physik (1–3 S.) + Reflexion (1–2 S.)
%%  Rohdaten, Auswerteformeln, ausgewertete Daten, Grafiken gehören zur Einlage.
%%  Lange Messwertlisten -> Anhang (4_messergebnis_anhang.tex), nicht in den Text.
%%%%%%%%%%%%%%%%%%%%%%%%%%%%%%%%%%%%%%%%%%%%%%%%%%%%%%%%%%%%%%%%%%%%%%%%%%%%%%%%
\section{Einlage: Eines Ihrer Ergebnisse (Messung oder Berechnung).}
\label{sec:messergebnis}

\begin{hinweise}
Hier nehmen Sie ein von Ihnen gemessenes Resultat (Rohdaten, Auswerteformeln und die ausgewerteten Daten und
Grafiken) und diskutieren es. Ein Bezug zu Ihrem Paper (Einlage~2) ist Voraussetzung.
\end{hinweise}

\subsection{Beschreibung}
\begin{hinweise}
Die Beschreibung gibt alle notwendigen Informationen an, die die Einlage beschreiben, also zum Beispiel das
Datum, die Art der Messung und eine kurze Inhaltsübersicht, aber keine Nacherzählung (min. $1/4$ Seite bis
max. $1/2$ Seite!).
\end{hinweise}
\begin{umfang}{Einlage 4 (Messergebnis): Beschreibung}{0.25}{0.5}
\platzhalter{Datum, Art der Messung, Messaufbau in Kürze -- laut Arbeitsplan: 3D-gedruckter Wackelstein
mit eingebautem Inertialsensor (Gyroskop) auf nivellierter Glasplatte (30\,cm\,$\times$\,30\,cm),
reproduzierbarer Anstoß mit definiertem Anfangs-$\omega$; eine der Messreihen 1--3 (Anfangs-$\omega$
variiert, Nick- vs.\ Roll-Anregung, $\delta$ variiert); Auswertung mit der eigenen Python-Datenpipeline;
Rohdaten in Anhang~\ref{anh:messdaten}, Ergebnis in Abbildung~\ref{fig:messergebnis}}
\end{umfang}

\subsection{Diskussion der Physik}
\begin{hinweise}
Beschreiben Sie hier die Physik dieser Einlage.
\begin{itemize}\itemsep0pt
\item Was haben Sie gemessen oder berechnet?
\item Was ist die Hauptaussage?
\item Was sind die Voraussetzungen für Ihre Messung?
\item Welches sind Ihre Auswerteformeln?
\item Was trägt ihre \textbf{wissenschaftliche Arbeit} zur Argumentation bei?
\item Welche Physik dieser Einlage ist bei Ihrem Schwerpunktthema anwendbar?
\item Wo gab es Schwierigkeiten beim Verständnis der Physik? Welche?
\item Wie haben Sie Ihr Paper aus der Einlage \flqq Eine wissenschaftliche Arbeit\frqq{} für diese Einlage verwendet? Warum? Wo?
\item Stimmen Ihre Ergebnisse mit den Aussagen aus dem Paper überein?
\end{itemize}
\textbf{Die Diskussion der Physik sollte mindestens eine Seite und maximal 3 Seiten lang sein. Die
Längenangaben sind ohne Titel, Bilder, Tabellen und Tabellen- oder Bildlegenden und ohne übergrosse Abstände
zu Gleichungen zu verstehen. Ich werde gegebenenfalls Ihre Auswertung anhand der beigelegten Daten
überprüfen.}

\smallskip Regeln für Abbildungen/Tabellen/Gleichungen (Limmer): jede Abbildung mit Bildunterschrift und im
Text referenziert; Schrift in Diagrammen $\geq 10$\,pt in Endgröße; Achsen \enquote{Größe / Einheit}
(z.\,B. $\omega_z$\,/\,rad\,s$^{-1}$, nicht $\omega_z$ [rad/s]); bei guter Übereinstimmung
Messung--Theorie/Simulation zusätzlich Residuenplot; Tabellen mit Überschrift; Messdaten nicht (nur) als
Tabelle, sondern (auch) als Diagramm; wichtige Gleichungen nummeriert, alle Symbole definiert.
Die Gerüste unten (Gleichung, Abbildung, Tabelle) sind auf das Projekt zugeschnitten -- anpassen oder löschen.
Bezug zum Paper ist Pflicht: Messergebnis mit der Umkehrzeit-Formel bzw.\ den Simulationen aus
\cite{\PaperKey} und mit der eigenen Simulation (Arbeitspaket 1.6) vergleichen.
\end{hinweise}
\begin{umfang}{Einlage 4 (Messergebnis): Diskussion der Physik}{1}{3}
\platzhalter{Was wurde gemessen, Hauptaussage, Voraussetzungen (Rollen ohne Gleiten auf der Glasplatte,
Kalibrierung des Gyroskops, \dots)}

% ---------- Gerüst: Auswerteformel mit Nummer und Verweis ----------
% Wachstum der Nick- bzw. Rollschwingung in der instabilen Drehrichtung (Arbeitspaket 5.2):
\begin{equation}
  A(t) = A_0\,\mathrm{e}^{\gamma t},
  \label{eq:wachstum}
\end{equation}
wobei $A$ die Amplitude der Nick- bzw.\ Rollschwingung, $A_0$ ihr Anfangswert und $\gamma$ die
Wachstumsrate der Instabilität ist. \platzhalter{Gleichung prüfen/ersetzen; Umkehrzeit $t_r$ aus den
Gyroskopdaten definieren (z.\,B.\ Zeit bis zum Vorzeichenwechsel von $\omega_z$); Bezug zum Paper:
\enquote{Nach \cite[Gl.~(\dots)]{\PaperKey} gilt für $t_r$ \dots}; Gl.~\eqref{eq:wachstum} im Text verwenden}

% ---------- Gerüst: Abbildung (Datei nach abbildungen/ legen) ----------
\begin{figure}[htbp]
  \centering
  %\includegraphics[width=0.8\textwidth]{messergebnis.pdf}
  \fbox{\parbox[c][4cm][c]{0.7\textwidth}{\centering\platzhalter{Abbildung: abbildungen/messergebnis.pdf,
    z.\,B.\ $\omega_z(t)$ aus dem Gyroskop mit Spin-Umkehr}}}
  \caption{\platzhalter{Aussagekräftige Bildunterschrift: Winkelgeschwindigkeit $\omega_z$ um die
    Hochachse über der Zeit $t$ für \dots, Messbedingungen, Ergebnis der Anpassung mit
    Gl.~\eqref{eq:wachstum}. Fremde Abbildungen: \cite[Abb.~\dots]{\PaperKey}}}
  \label{fig:messergebnis}
\end{figure}

% ---------- Gerüst: kurze Ergebnistabelle (Rohdaten in den Anhang!) ----------
\begin{table}[htbp]
  \centering
  \caption{\platzhalter{Tabellenüberschrift, z.\,B.\ gemessene Umkehrzeit $t_r$ in Abhängigkeit
    vom Anfangs-$\omega_0$ (Messreihe 1) im Vergleich zu Simulation und Paper}}
  \label{tab:ergebnis}
  \begin{tabular}{cccc}
    \toprule
    {$\omega_0$ / \si{\radian\per\second}} & {$t_r$ (Messung) / \si{\second}}
      & {$t_r$ (Simulation) / \si{\second}} & {$t_r$ (Paper) / \si{\second}} \\
    \midrule
    \dots & \dots & \dots & \dots \\
    \bottomrule
  \end{tabular}
\end{table}

\platzhalter{Diskussion: Abbildung~\ref{fig:messergebnis} und Tabelle~\ref{tab:ergebnis} im Text
verwenden, Unsicherheiten (Fehlerrechnung, Arbeitspaket 5.3), Vergleich Messung -- eigene Simulation --
Paper (5.4), Anwendbarkeit auf \Schwerpunkt}
\end{umfang}

\subsection{Reflexion}
\begin{hinweise}
In der Reflexion sollen Sie sich Gedanken machen über die Wirkung dieser Einlage und der Arbeit mit ihr auf
Ihren Lernfortschritt.
\begin{itemize}\itemsep0pt
\item In welchem Zusammenhang steht die Einlage mit Ihrem Schwerpunktthema?
\item Was haben Sie bei der Erstellung dieser Einlage gelernt?
\item Wie hat Ihnen die Einlage beim Verständnis Ihres Projektthemas geholfen?
\item Wo gab es Schwierigkeiten bei der Beschäftigung mit der Einlage?
\item Was hätten Sie anders gemacht, wenn Sie die Einlage nochmals erstellen müssten?
\item Was sagen Ihre Kommilitoninnen oder Kommilitonen dazu, als Kommentar oder wenn Sie mit Ihnen über die Einlage diskutieren?
\item Haben Sie in der Durchführung oder Auswertung der Messung logische oder physikalische Fehler entdeckt?
\item Gab es bei der Beschäftigung mit der Einlage eine überraschende neue Erkenntnis?
\end{itemize}
\textbf{Die Reflexion für eine Einlage sollte mindestens eine Seite und maximal 2 Seiten lang sein. Die
Längenangaben sind ohne Titel, Bilder, Tabellen und Tabellen- oder Bildlegenden zu verstehen.}
\end{hinweise}
\begin{umfang}{Einlage 4 (Messergebnis): Reflexion}{1}{2}
\platzhalter{Reflexion in Fließtext -- Leitfragen erkennbar beantworten}
\end{umfang}

%\subsection{Eigene Bewertungskriterien}
%...

%%%%%%%%%%%%%%%%%%%%%%%%%%%%%%%%%%%%%%%%%%%%%%%%%%%%%%%%%%%%%%%%%%%%%%%%%%%%%%%%
%%  EINLAGE 5: IHR VORTRAG ÜBER IHR PROJEKT
%%  Folien-Upload: 27.01.2027 | Vortrag: 28.01.2027 | vorgeschlagene Abgabe: 07.02.2027 (freiwillig)
%%  Struktur: Beschreibung (1/4–1/2 S.) + Reflexion (1–2 S.)
%%  Einlage = die Folien des Teils, den DU gehalten hast -> Anhang.
%%%%%%%%%%%%%%%%%%%%%%%%%%%%%%%%%%%%%%%%%%%%%%%%%%%%%%%%%%%%%%%%%%%%%%%%%%%%%%%%
\section{Einlage: Ihr Vortrag über ihr Projekt.}
\label{sec:vortrag}

\begin{hinweise}
Hier nehmen Sie die von Ihnen erstellten und präsentierten Folien des Projektvortrages und diskutieren sie
(nur den Teil, den Sie gehalten haben). Für die Moodle-Abgabe am 07.02.: \enquote{Die Folien ihres Vortrages
mit Kommentar und Reflexion dazu}.
\end{hinweise}

\subsection{Beschreibung}
\begin{hinweise}
Die Beschreibung gibt alle notwendigen Informationen an, die die Einlage beschreiben, also zum Beispiel das
Datum, die Art der Folien, die verwendete Software und eine kurze Inhaltsübersicht, aber keine Nacherzählung
(min. $1/4$ Seite, max. $1/2$ Seite!).
\end{hinweise}
\begin{umfang}{Einlage 5 (Vortrag): Beschreibung}{0.25}{0.5}
\platzhalter{Datum (\Vortragsdatum, Prüfung mit Anwesenheitspflicht), Anlass, Art der Folien und
Software (laut Arbeitsplan: LaTeX-Beamer mit der offiziellen Folien-Vorlage, Zitate auf jeder Folie),
Video der Spin-Umkehr (Arbeitspaket 3.2), Anzahl Folien, eigener Anteil vs.\ \Partner, kurze
Inhaltsübersicht; Probevortrag 26.\,01., Folien-Upload 27.\,01.\,2027; Folien in Anhang~\ref{anh:vortrag}}
\end{umfang}

\subsection{Reflexion}
\begin{hinweise}
In der Reflexion sollen Sie sich Gedanken machen über die Wirkung dieser Einlage und der Arbeit mit ihr auf
Ihren Lernfortschritt.
\begin{itemize}\itemsep0pt
\item In welchem Zusammenhang steht die Einlage mit Ihrem Projektthema?
\item In welchem Zusammenhang steht die Einlage mit Ihrem Schwerpunktthema?
\item Was haben Sie bei der Erstellung dieser Einlage gelernt?
\item Wie bewerten sie die Zusammenarbeit in Ihrer Gruppe?
\item Wie hat Ihnen die Einlage beim Verständnis Ihres Projektthemas geholfen?
\item Wie hat ihre \textbf{wissenschaftliche Arbeit} zur Argumentation beigetragen?
\item Wie hat Ihnen die Einlage beim Verständnis Ihres Schwerpunktthemas geholfen?
\item Wo gab es Schwierigkeiten bei der Beschäftigung mit der Einlage?
\item Was hätten Sie anders gemacht, wenn Sie die Einlage nochmals erstellen müssten?
\item Was sagen Ihre Kommilitoninnen oder Kommilitonen dazu, als Kommentar oder wenn Sie mit Ihnen über die Einlage diskutieren?
\item Enthält Ihre Präsentation logische oder physikalische Fehler?
\item Gab es bei der Beschäftigung mit der Einlage eine neue Erkenntnis?
\end{itemize}
\textbf{Die Reflexion für eine Einlage sollte mindestens eine Seite und maximal 2 Seiten lang sein. Die
Längenangaben sind ohne Titel, Bilder, Tabellen und Tabellen- oder Bildlegenden zu verstehen.}
Tipp: Fragen/Feedback aus der Diskussion nach dem Vortrag notieren und zitieren.
\end{hinweise}
\begin{umfang}{Einlage 5 (Vortrag): Reflexion}{1}{2}
\platzhalter{Reflexion in Fließtext -- Leitfragen erkennbar beantworten; einzelne Folien konkret
ansprechen (\enquote{Folie 4: \dots})}
\end{umfang}

%\subsection{Eigene Bewertungskriterien}
%...

%%%%%%%%%%%%%%%%%%%%%%%%%%%%%%%%%%%%%%%%%%%%%%%%%%%%%%%%%%%%%%%%%%%%%%%%%%%%%%%%
%%  EINLAGE 6: QUANTITATIVE AUSSAGE AUS EINEM ANDEREN VORTRAG
%%  vorgeschlagene Abgabe: 21.02.2027 (freiwillig)
%%  Struktur: Beschreibung (1/4–1/2 S.) + Diskussion der quantitativen Physik (1–3 S.) + Reflexion (1–2 S.)
%%  Vortrag/Folien der anderen Gruppe in literatur.bib eintragen (Vorlage VORLAGE-Vortrag).
%%%%%%%%%%%%%%%%%%%%%%%%%%%%%%%%%%%%%%%%%%%%%%%%%%%%%%%%%%%%%%%%%%%%%%%%%%%%%%%%
% \FremdVortragKey wird in einstellungen.tex gesetzt

\section{Einlage: Eine quantitative wissenschaftliche Aussage aus einem anderen Vortrag, die für Ihr Projekt eine Relevanz hat oder anwendbar ist.}
\label{sec:fremdvortrag}

\begin{hinweise}
Die Quelle dieser Einlage oder die Art der Einlage ist aus den von Ihnen gehörten Vorträgen (ausser dem
eigenen) frei wählbar. Sie muss mit Ihrem Schwerpunktthema einen Zusammenhang haben.

\smallskip\emph{Kandidaten aus der Themenliste WS 26/27 mit Bezug zu \Schwerpunkt:}
\enquote{Stabilität von rollenden Fahrrädern} (Stabilität aus linearisierten Bewegungsgleichungen,
Eigenwerte), \enquote{Untersuchung eines Doppelpendels} (numerische Integration, Chaos -- vgl.\ chaotisches
Verhalten bei schneller Drehung in \cite{\PaperKey}), \enquote{Dynamik eines rotierenden Tischtennisballs}
(Kreiseldynamik). Erst nach den Vorträgen am 28.\,01.\,2027 festlegen.
\end{hinweise}

\subsection{Beschreibung}
\begin{hinweise}
Die Beschreibung gibt alle notwendigen Informationen an, die die Einlage beschreiben, also zum Beispiel die
Gruppe, die den Vortrag gehalten hat, das Thema dieser Gruppe und eine kurze Inhaltsübersicht, aber keine
Nacherzählung (min. $1/4$ Seite, max. $1/2$ Seite!).
\end{hinweise}
\begin{umfang}{Einlage 6 (fremder Vortrag): Beschreibung}{0.25}{0.5}
\platzhalter{Gruppe, Thema, Datum, gewählte Aussage (Folie \dots) \cite{\FremdVortragKey}; kurze
Inhaltsübersicht; Folie(n) in Anhang~\ref{anh:fremdvortrag}}
\end{umfang}

\subsection{Diskussion der quantitativen Physik (oder der Mathematik)}
\begin{hinweise}
Wenn Sie für diese Einlage eine mathematische oder physikalische Aussage gewählt haben, beschreiben Sie hier
die Physik oder Mathematik dieser Einlage.
\begin{itemize}\itemsep0pt
\item Was ist die Hauptaussage?
\item Was sind die Voraussetzungen?
\item Wie hilft physikalisch die Physik (Mathematik oder Wissenschaft) dieser Einlage bei Ihrem Schwerpunktthema?
\item Welche quantitativen Aussagen unterlegt mit physikalischen Gleichungen sind für Sie besonders wichtig?
\item Wo gab es Schwierigkeiten beim Verständnis der Physik (Mathematik oder Wissenschaft)? Welche?
\item Wie können Sie mit dieser Aussage Ihr Projekt und Ihre Behandlung des Schwerpunktthemas verbessern oder korrigieren? Welchen Einfluss hätte dies?
\end{itemize}
\textbf{Die Diskussion der Physik sollte mindestens eine Seite und maximal 3 Seiten lang sein. Die
Längenangaben sind ohne Titel, Bilder, Tabellen und Tabellen- oder Bildlegenden und ohne übergrosse Abstände
zu Gleichungen zu verstehen.}
\end{hinweise}
\begin{umfang}{Einlage 6 (fremder Vortrag): Diskussion der Physik}{1}{3}
\platzhalter{Hauptaussage, Voraussetzungen, Gleichungen (nummeriert), Anwendung auf \Schwerpunkt,
quantitative Abschätzung des Einflusses auf das eigene Projekt}
\end{umfang}

\subsection{Reflexion}
\begin{hinweise}
In der Reflexion sollen Sie sich Gedanken machen über die Wirkung dieser Einlage und der Arbeit mit ihr auf
Ihren Lernfortschritt.
\begin{itemize}\itemsep0pt
\item In welchem Zusammenhang steht die Einlage mit Ihrem Projektthema?
\item In welchem Zusammenhang steht die Einlage mit Ihrem Schwerpunktthema?
\item Was haben Sie bei der Beschäftigung mit dieser Einlage gelernt?
\item Wie hat Ihnen die Einlage beim Verständnis Ihres Schwerpunktthemas geholfen?
\item Wo gab es Schwierigkeiten bei der Beschäftigung mit der Einlage?
\item Was hätten Sie anders gemacht, wenn Sie die Einlage erstellt hätten?
\item Was sagen Ihre Kommilitoninnen oder Kommilitonen dazu, als Kommentar oder wenn Sie mit Ihnen über die Einlage diskutieren?
\item Enthält die die verwendete Aussage logische oder physikalische Fehler?
\item Gab es bei der Beschäftigung mit der Einlage eine überraschende neue Erkenntnis?
\end{itemize}
\textbf{Die Reflexion für eine Einlage sollte mindestens eine Seite und maximal 2 Seiten lang sein. Die
Längenangaben sind ohne Titel, Bilder, Tabellen und Tabellen- oder Bildlegenden zu verstehen.}
\end{hinweise}
\begin{umfang}{Einlage 6 (fremder Vortrag): Reflexion}{1}{2}
\platzhalter{Reflexion in Fließtext -- Leitfragen erkennbar beantworten}
\end{umfang}

%\subsection{Eigene Bewertungskriterien}
%...


%% Freiwillig: weitere Einlagen (Vorlage kopieren, dann hier einbinden)
%%%%%%%%%%%%%%%%%%%%%%%%%%%%%%%%%%%%%%%%%%%%%%%%%%%%%%%%%%%%%%%%%%%%%%%%%%%%%%%%%
%%  VORLAGE FÜR FREIWILLIGE WEITERE EINLAGEN
%%  Datei kopieren (z.B. 7_datenblatt.tex), in Endabgabe einbinden:
%%      \input{einlagen/7_datenblatt}
%%  "Sie können eine beliebige Anzahl (>0) weiterer Einlagen zu Ihrem
%%   Schwerpunktthema einreichen. Die Struktur ist die gleiche wie bei den
%%   sechs Pflichteinlagen."  Beispiele laut Leitfaden: Grafik, Buchkapitel,
%%   Handout, Gesprächsnotiz, Datenblatt, selbst angefertigtes Objekt,
%%   Video-/Tonaufzeichnung, klassisches Protokoll.
%%%%%%%%%%%%%%%%%%%%%%%%%%%%%%%%%%%%%%%%%%%%%%%%%%%%%%%%%%%%%%%%%%%%%%%%%%%%%%%%
\section{Einlage: \platzhalter{Titel der Einlage}}
\label{sec:zusatz1}

\subsection{Beschreibung}
\begin{hinweise}
Um welche Art von Einlage handelt es sich? Woher stammt diese Einlage, wie bin ich darauf aufmerksam
geworden? Welche sind die wesentlichen Aussagen dieser Einlage (eigene Worte)? ($1/4$ bis max.\ $1/2$ Seite)
\end{hinweise}
\begin{umfang}{Zusatzeinlage: Beschreibung}{0.25}{0.5}
\platzhalter{Beschreibung}
\end{umfang}

%% Nur falls physikalischer Inhalt (dann 1–3 Seiten):
%\subsection{Diskussion der Physik}
%\begin{umfang}{Zusatzeinlage: Diskussion der Physik}{1}{3}
%...
%\end{umfang}

\subsection{Reflexion}
\begin{hinweise}
Wie habe ich mit dieser Einlage gearbeitet? Warum habe ich sie ausgewählt? Bezug zum Schwerpunktthema?
Wichtigste Erkenntnisse jenseits des wissenschaftlichen Inhalts? Bekanntes vertieft, neu gelernt, neue Fragen?
Einfluss auf den weiteren Lernprozess -- welche Konsequenzen? Was sagen Kommilitoninnen/Kommilitonen dazu?
Unstimmigkeiten? Überraschende Erkenntnis? (1--2 Seiten)
\end{hinweise}
\begin{umfang}{Zusatzeinlage: Reflexion}{1}{2}
\platzhalter{Reflexion}
\end{umfang}


%-----------------------------------------------------------------------------
%%%%%%%%%%%%%%%%%%%%%%%%%%%%%%%%%%%%%%%%%%%%%%%%%%%%%%%%%%%%%%%%%%%%%%%%%%%%%%%%
%%  GESAMTREFLEXION  (nach dem Hauptteil schreiben)   Umfang: 3–6 Seiten
%%  Gewicht 20 % der Note.
%%%%%%%%%%%%%%%%%%%%%%%%%%%%%%%%%%%%%%%%%%%%%%%%%%%%%%%%%%%%%%%%%%%%%%%%%%%%%%%%
\chapter{Gesamtreflexion}
\label{chap:gesamtreflexion}

\begin{hinweise}
In der Gesamtreflexion \cite{DK2008} überdenken Sie nochmals die Vorlesung und Ihr Portfolio. Sie sollten als
Ausgangspunkt Ihrer Überlegungen Ihr am Anfang ausgefülltes Motivationsformular verwenden (Anhang). Sie sollten
auf folgendes eingehen:
\begin{itemize}\itemsep0pt
\item Welche Erkenntnisse haben Sie aus den Portfolioeinlagen zu Ihrem Projekt gewonnen?
\item Welche Erkenntnisse haben Sie zum Schwerpunktthema gewonnen?
\item Was haben Sie zu Ihrer Art zu Lernen gelernt? Was war förderlich oder hinderlich?
\item Wie hat die Portfolioarbeit Ihre Arbeitsorganisation beeinflusst?
\item Was haben Sie an Allgemeinwissen dazugewonnen?
\item Was haben Sie über sich selber gelernt?
\item Wie beurteilen Sie Ihre Interaktion mit den Kommilitoninnen und Kommilitonen bei der Portfolioerstellung?
\item Wie beurteilen Sie ihren Studien-/Lernerfolg? Wie lautet ihre mit kritischer Bewertung? Welche Schlüsse ziehen Sie für Ihre weitere wissenschaftliche Bildung?
\item Welche Konsequenzen ziehen Sie aus den neu gewonnenen Erkenntnissen?
\item Welche \textbf{neue individuell bedeutsame Fragestellung} (oder Fragestellungen) hat sich aus den neu gewonnenen Erkenntnissen ergeben?
\item Was sind Ihre nächsten Schritte zur Auseinandersetzung mit den neu gewonnenen Fragestellungen?
\item Was ist Ihr Fazit?
\end{itemize}
Bitte schreiben Sie Ihren Text so, dass die geforderten Aussagen identifizierbar sind! Antworten aus nur einem
Satz sind meistens schwer bis unverständlich.
\textbf{Die Gesamtreflexion sollte mindestens drei Seiten und maximal sechs Seiten lang sein. Die Längenangaben
sind ohne Titel, Bilder, Tabellen und Tabellen- oder Bildlegenden zu verstehen.}

\smallskip\emph{Bewertungskriterien zusätzlich beachten:} begründete Einschätzung, inwieweit die
Auseinandersetzung mit dem Schwerpunktthema gelungen ist; fördernde/hemmende Faktoren; was bei Projekt- und
Schwerpunktthema geholfen hat; Ideen zur Fortführung der Forschung und welche Schritte dafür nötig wären;
Konsequenzen für Studium/Beruf. Wird \cite{DK2008} im eigenen Text nicht zitiert, erscheint es auch nicht im
Literaturverzeichnis.
\end{hinweise}

\begin{umfang}{Gesamtreflexion}{3}{6}

\section*{Erkenntnisse zu Projekt- und Schwerpunktthema}
\platzhalter{Rückblick: physikalische Erkenntnisse aus den Einlagen, bezogen auf \Schwerpunkt{}
(z.\,B.\ Übereinstimmung von Messung, eigener Simulation und Garcia-Hubbard-Formel aus \cite{\PaperKey})}

\section*{Lernprozess, Arbeitsorganisation und Zusammenarbeit}
\platzhalter{Art zu lernen, förderlich/hinderlich, Arbeitsorganisation (Soll-Ist-Vergleich mit dem
Arbeitsplan, Anhang~\ref{anh:arbeitsplan}), über sich selbst, Interaktion mit \Partner{} und
Kommilitoninnen und Kommilitonen}

\section*{Bewertung des Lernerfolgs}
\platzhalter{kritische Bewertung, Vergleich mit dem Motivationsformular, Schlüsse für die weitere Bildung}

\section*{Neue Fragestellungen und nächste Schritte}
\platzhalter{mindestens eine neue individuell bedeutsame Fragestellung + begründete nächste Schritte}

\section*{Fazit}
\platzhalter{Fazit}

\end{umfang}


%-----------------------------------------------------------------------------
\printbibliography[heading=bibintoc,title={Referenzen und Literaturverzeichnis}]
\begin{hinweise}
Ein Literaturverzeichnis enthält auch Hinweise auf Hilfen, Texte, Grafiken, Aussagen von anderen Menschen,
z.B. von Mitstudierenden! (Alles was Sie übernommen haben!) In dieses Verzeichnis gehören neben korrekt
zitierten Literaturstellen (nicht die Webadresse) auch alle anderen Verweise wie auf Aussagen von anderen
Personen, Webstellen, die verwendet wurden, und was immer sonst nicht auf Ihrem Mist gewachsen ist.
Mit Jabref oder Citavi können Sie Ihre Literaturliste pflegen!
\end{hinweise}

%-----------------------------------------------------------------------------
\appendix
\chapter*{Anh\"{a}nge}
\addcontentsline{toc}{chapter}{Anh\"{a}nge}
\hinweis{Hier kommen Ihre Anhänge hin. Das ist das Motivationsformular, alle Einlagen zu den sechs
Pflichtabgaben sowie weitere Materialien nach Ihrer Wahl. Die einzelnen Anhänge sollten mit
\textbackslash chapter\{welcher Anhang\} betitelt werden!}

\chapter{Motivationsformular}
\label{anh:motivationsformular}
%%%%%%%%%%%%%%%%%%%%%%%%%%%%%%%%%%%%%%%%%%%%%%%%%%%%%%%%%%%%%%%%%%%%%%%%%%%%%%%%
%%  MOTIVATIONSFORMULAR – Fragen wörtlich aus
%%  "Motivationsformular für Portfolio-2026-2027-tex.tex" (Karagiannakis/Marti).
%%  Deine Antworten in \antwort{...} eintragen. Wird verwendet von
%%   - Manevski_Motivationsformular.tex   (Abgabe 21.10.2026)
%%   - Endabgabe (Anhang "Motivationsformular")
%%  Später NICHT mehr inhaltlich ändern: es dokumentiert deinen Stand zu Beginn.
%%  -> Antworten aus der am 21.10.2026 in Moodle abgegebenen Fassung hier 1:1 einfügen
%%     (der Text liegt nicht im Projektordner).
%%%%%%%%%%%%%%%%%%%%%%%%%%%%%%%%%%%%%%%%%%%%%%%%%%%%%%%%%%%%%%%%%%%%%%%%%%%%%%%%

Die folgenden sechs Fragen sollen Ihnen helfen, Ihre Motivation und Ihre Vorkenntnisse zu bewerten.
Sie k\"{o}nnen Teile der Gesamtreflexion wie auch der Einleitung mit Hilfe Ihrer Antworten formulieren.

\begin{enumerate}

\item Ich nehme am Projektpraktikum teil, weil \ldots

\antwort{\platzhalter{Antwort 1}}

\item Ich bringe folgende mathematische, physikalische und sonstige Vorkenntnisse und / oder
praktische Erfahrungen mit \ldots

\antwort{\platzhalter{Antwort 2}}

\item Das Projektpraktikum wird f\"{u}r mich erfolgreich sein, wenn Folgendes eintritt \ldots

\antwort{\platzhalter{Antwort 3}}

\item Folgendes darf hier nicht passieren \ldots

\antwort{\platzhalter{Antwort 4}}

\item Das ist mir sonst noch wichtig \ldots

\antwort{\platzhalter{Antwort 5}}

\item Ich bin bereit, Verantwortung sowohl f\"{u}r das Gelingen des Projektpraktikums als auch f\"{u}r
meinen eigenen Lernprozess zu \"{u}bernehmen und werde folgenden Beitrag leisten \ldots

\antwort{\platzhalter{Antwort 6}}

\end{enumerate}

Quelle: E. Karagiannakis, Universit\"{a}t Ulm, Portfolio-Weiterbildung 2012


%%%%%%%%%%%%%%%%%%%%%%%%%%%%%%%%%%%%%%%%%%%%%%%%%%%%%%%%%%%%%%%%%%%%%%%%%%%%%%%%
%%  ANHANG zu Einlage 1: der Arbeitsplan selbst
%%  Variante A (aktiv): Excel-Export Arbeitsplan_Wackelstein.pdf (A3 quer, 2 Seiten)
%%              = Kopie aus Arbeitsplan/Aktueller Arbeitsplan/ -> anhaenge/arbeitsplan.pdf.
%%              Nach jeder Änderung am Excel-Plan neu als PDF exportieren und ersetzen.
%%  Variante B: Gantt-Diagramm direkt in LaTeX (pgfgantt) mit den Daten aus der
%%              Excel-Datei (Stand 22.10.2026). Zum Aktivieren Variante A auskommentieren
%%              und \iffalse/\fi um Variante B entfernen.
%%%%%%%%%%%%%%%%%%%%%%%%%%%%%%%%%%%%%%%%%%%%%%%%%%%%%%%%%%%%%%%%%%%%%%%%%%%%%%%%

%% ---- Variante A ----
\anhangpdf[pages=-,landscape]{Arbeitsplan (Einlage 1)}{anh:arbeitsplan}{anhaenge/arbeitsplan.pdf}

%% ---- Variante B ----
\iffalse
\chapter{Arbeitsplan (Einlage 1)}
\label{anh:arbeitsplan}

\hinweis{Daten aus Arbeitsplan\_Wackelstein.xlsx (Stand 22.\,10.\,2026), auf die oberen Ebenen
gekürzt. Ein Tag = 0{,}58\,mm. K = Kristijan, J = Janosch, B = beide.}

\begin{figure}[htbp]
\centering
\begin{ganttchart}[
    time slot format=isodate,
    x unit=0.58mm, y unit chart=0.42cm, y unit title=0.5cm,
    hgrid, vgrid={*6{draw=none},dotted},
    title label font=\scriptsize, bar label font=\tiny,
    group label font=\scriptsize\bfseries, milestone label font=\tiny\itshape,
    bar height=0.55, group height=0.35,
    bar/.append style={fill=blue!25},
    group/.append style={fill=black!60},
    milestone/.append style={fill=red!70},
  ]{2026-10-12}{2027-03-21}
  \gantttitlecalendar{year, month=shortname} \\
  \ganttgroup{1 Theorie \& Literatur}{2026-10-15}{2026-12-17} \\
  \ganttbar{1.1 Literaturrecherche (B)}{2026-10-15}{2026-11-04} \\
  \ganttbar{1.2/1.3 Paper wählen + durcharbeiten (B)}{2026-10-29}{2026-11-10} \\
  \ganttbar[bar/.append style={fill=orange!50}]{1.4 Bewegungsgleichungen (K)}{2026-10-29}{2026-11-25} \\
  \ganttbar{1.5 Trägheitstensor aus STL (J)}{2026-10-29}{2026-11-25} \\
  \ganttbar[bar/.append style={fill=orange!50}]{1.6 Simulation Umkehrdynamik (K)}{2026-11-19}{2026-12-17} \\
  \ganttbar{1.7 Vorhersage Umkehrzeit (J)}{2026-11-19}{2026-12-17} \\
  \ganttgroup{2 Materialien}{2026-10-12}{2026-11-26} \\
  \ganttbar{2.1.1 Sensor und Akku festlegen (B)}{2026-10-12}{2026-10-19} \\
  \ganttbar{2.1.2 Auftrag E-Werkstatt (B)}{2026-10-22}{2026-11-09} \\
  \ganttbar[bar/.append style={fill=orange!50}]{2.1.3/2.1.4 Firmware + Kalibrierung (K)}{2026-11-05}{2026-11-26} \\
  \ganttbar{2.2.1 Prototyp Stein, 3D-Druck (J/B)}{2026-10-12}{2026-11-08} \\
  \ganttbar{2.2.2 Messkörper mit Sensorfach (J)}{2026-11-05}{2026-11-12} \\
  \ganttgroup{3 Versuchsaufbau}{2026-10-22}{2026-12-14} \\
  \ganttbar{3.1 Glasplatte besorgen, testen, nivellieren}{2026-10-22}{2026-10-29} \\
  \ganttbar[bar/.append style={fill=orange!50}]{3.2 Video (K)}{2026-10-29}{2026-12-14} \\
  \ganttbar{3.3 Reproduzierbarer Anstoß}{2026-10-29}{2026-12-14} \\
  \ganttgroup{4 Messungen}{2026-10-15}{2027-02-11} \\
  \ganttbar{4.1 Laborjournal (B)}{2026-10-15}{2027-02-11} \\
  \ganttbar{4.2 Testmessung (B)}{2026-12-03}{2026-12-07} \\
  \ganttbar{4.3--4.5 Messreihen 1--3 (B)}{2026-12-07}{2027-01-21} \\
  \ganttbar{4.6 Nachmessungen / Puffer (B)}{2027-01-21}{2027-02-11} \\
  \ganttgroup{5 Auswertung}{2026-11-26}{2027-01-19} \\
  \ganttbar[bar/.append style={fill=orange!50}]{5.1 Datenpipeline Python (K)}{2026-11-26}{2026-12-17} \\
  \ganttbar{5.2--5.5 Fits, Fehlerrechnung, Vergleich (B)}{2026-12-07}{2027-01-19} \\
  \ganttgroup{6 Vortrag}{2027-01-14}{2027-01-28} \\
  \ganttbar{6.1/6.2 Folien + Probevortrag (B)}{2027-01-14}{2027-01-26} \\
  \ganttgroup{7 Portfolio (individuell)}{2026-10-15}{2027-03-15} \\
  \ganttbar{Einlagen schreiben}{2026-10-22}{2027-02-20} \\
  \ganttbar{Einleitung + Gesamtreflexion}{2027-02-22}{2027-03-07} \\
  \ganttbar{Korrektur + Puffer}{2027-03-08}{2027-03-13} \\
  \ganttmilestone{Arbeitsplan}{2026-10-28} \\
  \ganttmilestone{Wiss.\ Arbeit}{2026-11-11} \\
  \ganttmilestone{Protokoll}{2026-12-09} \\
  \ganttmilestone{Messergebnis}{2027-01-20} \\
  \ganttmilestone{Vortrag (Prüfung)}{2027-01-28} \\
  \ganttmilestone{Endabgabe Portfolio}{2027-03-15}
\end{ganttchart}
\caption{Arbeitsplan von \IhrName{} und \Partner{} für das Projekt \enquote{\Projekt}
(gemeinsam erstellt, Stand 22.\,10.\,2026). Orange: eigene Arbeitspakete, rote Rauten: Abgabetermine.}
\label{fig:arbeitsplan}
\end{figure}
\fi

%% ---- KI-Nutzung (Pflicht laut Zitierregeln): alle Prompts + ausgefüllter Fragebogen ----
%% erzeugt aus Fragebogen_Arbeitsplan.md; bei Änderungen am Fragebogen neu erzeugen
%% AUTOMATISCH ERZEUGT (make_ki_anhang.py) -- Prompts und Inputs der KI-Nutzung zu Einlage 1
\chapter{KI-Nutzung zu Einlage 1: Prompts und Inputs}
\label{anh:ki-arbeitsplan}

Beschreibung und Reflexion der Einlage Arbeitsplan (Abschnitt~\ref{sec:arbeitsplan}) sind von einem
KI-System formuliert worden \cite{Manevski2026KIArbeitsplan}. Verwendetes System: Claude Opus 5.5
(Anthropic) in Claude Code, 28.\,10.\,2026. Weitere Inputs: der gemeinsame Arbeitsplan
(Arbeitsplan\_Wackelstein.xlsx), die Kursunterlagen im Projektordner und die Arbeitsplan-Abgabe von
Janosch Lustig, die nur als Vergleich für Umfang und Stil diente. Im Folgenden stehen alle meine Prompts
in der Reihenfolge, in der ich sie eingegeben habe, und danach der von der KI erstellte und von mir
ausgefüllte Fragebogen.

\section*{Prompts}
{\small

\par\medskip\noindent\textbf{Prompt 1:}\\
{\ttfamily\raggedright
schau dir Moritz Beispielportfolio an (C:\textbackslash{}Users\textbackslash{}KIKO\textbackslash{}Documents\textbackslash{}!UNI\textbackslash{}Projekt-Praktikum\textbackslash{}Formulare für Abgaben\textbackslash{}Moritz Beispiel)\\
\smallskip
passe alle datei darin an sodass es für mich zugeschnitten ist und alle informatione beinhaltet die im projekt zu finden sind\\
speichere meine angepasste Kopie im ''Formulare für Abgaben'' Ordner als Kristijan Protfolio\\
}
\par\medskip\noindent\textbf{Prompt 2:}\\
{\ttfamily\raggedright
Nicht jetzt\\
}
\par\medskip\noindent\textbf{Prompt 3:}\\
{\ttfamily\raggedright
ok jetzt eine neue aufgabe\\
findest du in dem projekt eine datei die beschreibt wie die Reflexion des Arbeitsplans aussehen soll (inhalt, länge, etc)?\\
}
\par\medskip\noindent\textbf{Prompt 4:}\\
{\ttfamily\raggedright
{[}Eingefügter Text der Moodle-Aufgabe:{]}\\
Einlage: Arbeitsplan – Aufgabe\\
Termin: 28. Oktober 2026, 23:59 Uhr\\
\smallskip
1. Erstellen Sie gemeinsam mit Ihrer Partnerin / Ihrem Partner einen Arbeitsplan für Ihr Projekt. Besprechen Sie den Plan mit den Betreuern. Als Ideengeber können Sie den Arbeitsplan zum Mittagessen kochen nehmen. Es ist sinnvoll, den Plan im Gantt-Format zu erstellen. Neben spezialisierter Software können Sie auch Excel oder Libre Office Calc verwenden.\\
2. Schreiben Sie eine Reflexion über den Arbeitsplan. Eine Reflexion ist keine Nacherzählung. Verwenden Sie die zur Verfügung gestellten Template oder strukturieren Sie Ihre Abgabe danach.\\
3. Für diese Einlage verwenden alle Gruppenmitglieder den gleichen gemeinsamen Arbeitsplan, schreiben jedoch die Reflexionen einzeln und unabhängig voneinander.\\
4. Laden Sie den Arbeitsplan und die Reflexion hoch.\\
   * Akzeptierte Formate für alle Abgaben: Adobe PDF (Microsoft DOCX bitte zuert in PDF speichern). Wandeln Sie andere Formate in diese Formate um.\\
   * Umfang Beschreibung der Einlage: Mindestens: 1/4-Seite, Maximal 1/2 Seite\\
   * Umfang Reflexion: Mindestens 1 Seite, maximal 2 Seiten!\\
   * Alle Längenangaben sind ohne Titel, Abbildungen, Abbildungslegenden, Tabellen und Tabellenlegenden (oder Ähnliches) zu verstehen.\\
5. Vergessen Sie nicht nach dem Hochladen den Abgabeknopf zu drücken.\\
\smallskip
Dazugehörige Dateien finden Sie in Unterlagen zum Projektpraktikum, Einlage während Semester Template und Unterlagen zum Projektpraktikum, Arbeitsplan\\
}
\par\medskip\noindent\textbf{Prompt 5:}\\
{\ttfamily\raggedright
ok für den arbeitsplan\\
\smallskip
1. Einlage:\\
   1. was zur hölle ist eine ''Einlage''\\
   2. Stelle mir einen Fragebogen den ich beantworten soll damit du daraus die Einlage schreiben kannst\\
2. Reflexion:\\
   1. was ist mit ''Reflexion'' gemeint\\
   2. Stelle mir einen riesen fragebogen den ich beantworten soll damit du damit die Reflexion schreiben kannst\\
}
\par\medskip\noindent\textbf{Prompt 6:}\\
{\ttfamily\raggedright
{[}Anhang: pp\_reflexion\_arbeitsplan (1).pdf – Arbeitsplan-Abgabe von Janosch Lustig{]}\\
die .md file habe ich nun ausgefüllt, lesen sie nochmals durch und mache damit den text.\\
\smallskip
hier ist noch die in der .md angesprochene version von janosch\\
\smallskip
übertreib es nicht mit der länge\\
}
\par\medskip\noindent\textbf{Prompt 7:}\\
{\ttfamily\raggedright
Mit KI-Zitat schreiben\\
}
}

\section*{Fragebogen mit meinen Antworten}
{\small\raggedright

\par\noindent Antworte direkt unter jede Frage, \textbf{Stichpunkte reichen}. Ehrlich und konkret ist besser als schön:\par
\par\noindent „Ich hab die E-Werkstatt-Wartezeit unterschätzt, weil …“ ist Gold, „Planung ist wichtig“ ist wertlos.\par
\par\noindent Fragen, zu denen dir nichts einfällt: leer lassen oder „weiß nicht“. Nichts wird erfunden.\par
\subsection*{Teil 1 – Beschreibung der Einlage (1/4 bis 1/2 Seite)}
\par\noindent Ziel: Ein Leser soll wissen, \emph{was} der Arbeitsplan ist, \emph{woher} er kommt und \emph{was drinsteht} –\par
\par\noindent ohne ihn nachzuerzählen.\par
\par\noindent \textbf{1.1} Womit habt ihr den Plan gebaut? (Excel? Version? Mit Vorlage oder komplett selbst?)\par
\par\noindent\emph{Antwort:} Wir haben den Plan in Excel Gebaut. Haben uns an einer Baumstruktur pro parallel laufendem thema ein beispiel genommen.\par\smallskip
\par\noindent \textbf{1.2} Wann genau habt ihr ihn erstellt? (erste Version an welchem Tag, letzte Änderung? Der Plan sagt „Stand 22.10.2026“ – stimmt das noch?)\par
\par\noindent\emph{Antwort:} Der Plan wurde seit dem 22.10.2026 nicht verändert und wurde an genau diesem tag fertiggestellt und mit Prof. Marti durchgesprochen\par\smallskip
\par\noindent \textbf{1.3} Wer hat was daran gemacht? (Struktur, Aufgabenliste, Formatierung/Formeln, Termine-Blatt …)\par
\par\noindent\emph{Antwort:} Komplette zusammenarbeit in jedem bereich des plans\par\smallskip
\par\noindent \textbf{1.4} Wie lange habt ihr ungefähr daran gearbeitet (Stunden, wie viele Sitzungen)?\par
\par\noindent\emph{Antwort:} nicht getrackt\par\smallskip
\par\noindent \textbf{1.5} Was habt ihr als Vorlage / Ideengeber benutzt? (Beispiel „Mittagessen kochen“ aus Moodle? Petros' Plan – was genau habt ihr davon übernommen: Struktur, Termine, Idee mit Python/CSV?)\par
\par\noindent\emph{Antwort:} Mittagessen kochen als beispiel für die detalierten zweige des baums\par\smallskip
\par\noindent \textbf{1.6} Habt ihr den Plan mit den Betreuern besprochen? Wenn ja: wann, mit wem (Name), und was haben sie gesagt?\par
\par\noindent\emph{Antwort:} ja wie gesagt mit prof. marti und dieser hat ein paar fragen gestellt die wir im beantworten konnten. dannach hat er abgenickt\par\smallskip
\par\noindent \textbf{1.7} Was ist für dich die wichtigste Aussage des Plans in einem Satz? (z. B. „Kritischer Pfad ist der Sensor“, „Messungen erst ab Dezember möglich“ …)\par
\par\noindent\emph{Antwort:} für mich ist die wichtigeste aussage des plans die super leicht zu erkennenden aufgaben die parallel laufen und zusätzlich auch noch zusammenhängen damit man möglichst effizient ist\par\smallskip
\par\noindent \textbf{1.8} Welche Phasen / Meilensteine sind aus deiner Sicht die entscheidenden?\par
\par\noindent\emph{Antwort:} erster meilenstein wäre der fertige prototyp (das er funktioniert und man damit anfangen kann zu messen), zweiter meilenstein wäre die funktionierende simulation sodass man seine ergebnisse direkt validieren/vergleichen kann, dritter meilenstein ist dann eine schöne saubere messreihe womit man ordentlich arbeiten kann\par\smallskip
\par\noindent \textbf{1.9} Gibt es Besonderheiten am Plan, die erwähnt werden sollten? (Farbcodierung nach Person, automatische Balken, Statusspalte, Termine-Blatt mit Quellen, Praktikumspause …)\par
\par\noindent\emph{Antwort:} alles ist nach person aufgeteilt, die statusspalte ist super zum nachverfolgen was der andere macht, die prozentangabe des themenblocks ist praktisch um sich zu motivieren oder nichts zu vergessen und alle abgaben sind eingetragen\par\smallskip
\par\noindent \textbf{1.10} Ist der Plan als „lebendes Dokument“ gedacht? Wie oft wollt ihr ihn aktualisieren?\par
\par\noindent\emph{Antwort:} der plan darf nicht verändert werden. er wird von uns nur dazu genutzt um unseren fortschritt zu tracken.\par\smallskip
\subsection*{Teil 2 – Reflexion (1 bis 2 Seiten)}
\par\noindent Ziel: Nicht \emph{was} im Plan steht, sondern \emph{was das Erstellen mit dir gemacht hat} – was du gelernt,\par
\par\noindent gemerkt, falsch eingeschätzt hast und was du daraus folgerst.\par
\par\medskip\noindent\textbf{A – Wie lief das Erstellen?}\par
\par\noindent \textbf{A1} Wie seid ihr vorgegangen? (erst Liste, dann Zeiten? Rückwärts von den Abgaben geplant? Brainstorming?)\par
\par\noindent\emph{Antwort:} wir konnten anfangs nicht in worte fassen was wir uns vorgestellt haben wie der plan auszusehen hat. dann haben wir gebrainstormed und ''prototypen'' gezeichnet.\par\smallskip
\par\noindent \textbf{A2} Was war der erste Entwurf, und wie sehr hat sich der Plan bis zur jetzigen Version verändert? Was wurde umgeworfen und warum?\par
\par\noindent\emph{Antwort:} der plan ist nicht direkt in excel entstanden sondern im kopf dann mehrere itterationen auf dem blatt und dann erst als grober ansatz in excel\par\smallskip
\par\noindent \textbf{A3} Wie habt ihr entschieden, wer welche Aufgabe übernimmt? Nach Interesse, Können, Zeit?\par
\par\noindent\emph{Antwort:} unsere themenschwerpunkte waren uns klar. daher haben wir uns zu jedem schwerpunkt überlegt welche aufgaben alle zu erledigen sind. diese aufgaben haben wir dann alle runtergeschrieben. und untereinander aufgeteilt. dannach haben wir überlegt welche aufgaben gleichzeitig zusammenhängen und in welchem zeitabschnitt des PPs sie zu erledigen sind damit es am sinnvollsten ist. so haben wir dann die aufgaben gruppiert und aufgeteilt wann diese zu erledigen sind. jedoch können wir ja nicht jeder aufgabe gleich viel zeit geben, also haben wir noch disskutiert welche aufgabe wie viel zeit in anspruch nehmen wird, wobei die GP erfahrung der letzten 2 Jahre geholfen hat.\par\smallskip
\par\noindent \textbf{A4} Gab es Diskussionen oder Meinungsverschiedenheiten mit Janosch? Worüber? Wie gelöst?\par
\par\noindent\emph{Antwort:} es gab direkt am anfang meinungsverschiedenheiten die wir durch argumente und vielleicht vergessene themen alles schön lösen konnten sodass am ende beide glücklich sind und das beste aus beiden welten zusammenfindet\par\smallskip
\par\noindent \textbf{A5} Wie lief die Arbeit mit Excel selbst? (Formeln, bedingte Formatierung – leicht, nervig, neu für dich?)\par
\par\noindent\emph{Antwort:} frage ist todes unwichtig\par\smallskip
\par\medskip\noindent\textbf{B – Bezug zum Projektthema (Wackelstein)}\par
\par\noindent \textbf{B1} Was wusstest du über den Wackelstein / das Projekt \emph{bevor} ihr den Plan gemacht habt, und was danach?\par
\par\noindent\emph{Antwort:} ich wusste davor nur das er existiert und grob wie seine funktionalität ist. nach dem erstellen des plans ist mir die aufteilung des verständnisses klarer geworden. wie man eine funktion des steins betrachten sollte damit man sie am ende intuitiver oder besser versteht und was sich alles in der theorie versteht an das man nicht sofort nur beim anschauen des steins erkennt, was jedoch klar wird wenn man sich ein ziel für den versuch überlegt und dann die aufgaben ausdenken muss die einen zum ziel führen.\par\smallskip
\par\noindent \textbf{B2} Welche Arbeitsschritte sind dir erst beim Planen bewusst geworden? (z. B. Glasplatte testen, Reibung, Nivellieren, reproduzierbarer Anstoß, Sensor kalibrieren)\par
\par\noindent\emph{Antwort:} reibung, reproduzierbarer anstoß, sonsorkalibrierung nie als konkreten schritt angesehen welcher aber nach dem daran denken unfassbar wichtig und praktisch ist\par\smallskip
\par\noindent \textbf{B3} Wo hat das Planen dir gezeigt, dass das Projekt aufwendiger / einfacher ist als gedacht?\par
\par\noindent\emph{Antwort:} bei der theorie und der vorbereitung (das 3D model muss für so einen niechigen effekt sozusagen PERFEKT sein)\par\smallskip
\par\noindent \textbf{B4} Welche Teile des Projekts hängen voneinander ab, und hast du das vorher gesehen? (Sensor $\rightarrow$ Messkörper $\rightarrow$ Messungen; Trägheitstensor $\rightarrow$ Simulation)\par
\par\noindent\emph{Antwort:} die messungen hängen von der qualität des körpers ab, was klar ist. jedoch war es vorerst unklar das die kleinste veränderung der startbedingungen eine so wichtige rolle spielen wird.\par\smallskip
\par\medskip\noindent\textbf{C – Bezug zu deinem Schwerpunkt (Bewegungsgleichungen + Simulation der Spin-Umkehr)}\par
\par\noindent \textbf{C1} Ist dieser Schwerpunkt-Titel richtig? Wenn nicht: wie heißt dein Schwerpunkt genau, und was ist der von Janosch?\par
\par\noindent\emph{Antwort:} das ist meiner ja\par\smallskip
\par\noindent \textbf{C2} Wo taucht dein Schwerpunkt im Plan auf (welche Zeilen), und wie viel Zeit hast du dafür eingeplant? Realistisch?\par
\par\noindent\emph{Antwort:} das sollst du rausfinden (claude), das steht im plan\par\smallskip
\par\noindent \textbf{C3} Was brauchst du von Janosch, damit du deinen Teil machen kannst (z. B. Trägheitstensor aus der STL), und was braucht er von dir?\par
\par\noindent\emph{Antwort:} alle wichtigen größen des körpers und seine eigenschaften.\par\smallskip
\par\noindent \textbf{C4} Hat dir der Plan geholfen zu verstehen, was dein Schwerpunkt konkret bedeutet? Was war vorher unklar?\par
\par\noindent\emph{Antwort:} ja hat er da ich die chance hatte mich genauer damit zu befassen wie ich meine aufgaben einzuteilen habe und womit sie direkt oder indirekt zusammenhängen.\par\smallskip
\par\noindent \textbf{C5} Warum hast du (bzw. habt ihr) gerade dir die Aufgaben Bewegungsgleichungen, Simulation, Firmware, Kalibrierung, Datenpipeline gegeben? Was reizt dich daran?\par
\par\noindent\emph{Antwort:} mich reizt die programmierungm, theorie und simulaton da man dabei ein tieferes verständnis bekommt da man sich um die kleinstens details kümmern muss an die man sonst, wenn man zb eine simulation nicht als ziel hätte, nie denken würde und die dann am ende im verständnis fehlen würden\par\smallskip
\par\medskip\noindent\textbf{D – Was hast du gelernt?}\par
\par\noindent \textbf{D1} Was hast du über Projektplanung allgemein gelernt? (Aufgaben zerlegen, Puffer, Parallelität, externe Abhängigkeiten …)\par
\par\noindent\emph{Antwort:} das puffer sehr wichtig ist wenn man einen gegebenen zeitraum und etwas unbekanntes vor sich hat, es kann ja schlieslich alles schief gehen. gelernt habe ich auch das es sehr wichtig ist zu itterieren, wenn man an einem schon bekanntem thema nach einiger zeit und damit mit einem andern winkel arbeitet erkannt man was man vergessen hat oder hinterfragt die richtigkeit und wichtigkeit des schon erarbeitetem.\par\smallskip
\par\noindent \textbf{D2} Was hast du über dich selbst gelernt? (Neigst du zum Unterschätzen? Planst du gern / ungern? Arbeitest du lieber spontan?)\par
\par\noindent\emph{Antwort:} ich habe über mich gelern das ich eher dazu neige die zeit die man für etwas hat zu unterschätzen. dessen bin ich mir jedoch schon bewusst wordurch ich dafür kompensieren kann, so wie in diese plan zu sehen ist.\par\smallskip
\par\noindent \textbf{D3} Hattest du vorher schon mal so einen Plan gemacht (Schule, Job, andere Praktika)? Unterschied?\par
\par\noindent\emph{Antwort:} keinen solch detalierten plan\par\smallskip
\par\noindent \textbf{D4} Welche Vorlesung / welches Modul hätte dir dabei geholfen oder hat geholfen?\par
\par\noindent\emph{Antwort:} dabei hat die erfahrung aus dem GP (Grundpraktikum) geholfen da ich dort die erfahrung gemacht habe das aufgaben bei denen man denkt das sie schnell erledigt sind eigentlich nie schnell erledigt sind, wodurch ich in diesem plan (hoffentlich) nicht unterschätzt habe\par\smallskip
\par\medskip\noindent\textbf{E – Schwierigkeiten}\par
\par\noindent \textbf{E1} Was war am schwersten? (Dauer schätzen, Reihenfolge, wissen was man noch nicht weiß, Excel, Absprache …)\par
\par\noindent\emph{Antwort:} ich würde nicht sagen das mir etwas sehr schwer fiel. am anfang war der gedanke einen solch detalierten plan zu erstellen einschüchternd, hat aber nach dem beginnen nur positives gebracht wodurch am ende eher zufriedenheit als frust oder schwierigkeit zu emfinden war.\par\smallskip
\par\noindent \textbf{E2} Welche Dauern konntest du gar nicht einschätzen, und wie habt ihr dann geschätzt?\par
\par\noindent\emph{Antwort:} da es schwer ist mit dem 3D druck und teilliefenrungen zu planen haben wir dafür ordentlich puffer eingebaut damit es an solchen banalen dingen nicht scheitert.\par\smallskip
\par\noindent \textbf{E3} Wo hängt ihr von anderen ab (E-Werkstatt, 3D-Druck, Betreuer, Material)? Wie seid ihr damit umgegangen?\par
\par\noindent\emph{Antwort:} von unserem betreuer da er unseren Stein 3D druckt und wir hoffen müssen das er dauerhaft erreichbar ist falls es änderungen oder sowas gibt\par\smallskip
\par\noindent \textbf{E4} Gab es Aufgaben, bei denen du beim Planen nicht wusstest, was sie eigentlich beinhalten?\par
\par\noindent\emph{Antwort:} nicht wirlich, und auch wenn das am anfang so gewesen sein sollte ist es noch WÄHREND dem planen klar geworden.\par\smallskip
\par\medskip\noindent\textbf{F – Feedback von anderen}\par
\par\noindent \textbf{F1} Was hat Janosch zum Plan bzw. zu deinen Ideen gesagt?\par
\par\noindent\emph{Antwort:} es gab kritik sowe positives, an genaue wortlaute kann ich mich nicht errinnern. am ende waren wir beide zufrieden, was am wichtigsten ist\par\smallskip
\par\noindent \textbf{F2} Was haben die Betreuer gesagt? (Name, Datum, Inhalt – wird zitiert!)\par
\par\noindent\emph{Antwort:} wir haben wie gesagt mit marit gesprochen, dieser hat den plan in der gezeigten version abgenickt nachdem er fragen gestellt hat um den plan genauer zu verstehen (kein zitat, das gespräch war mündlich am Donnerstag den 22.10.2026)\par\smallskip
\par\noindent \textbf{F3} Hast du den Plan Kommilitonen gezeigt (Petros oder andere)? Was kam als Kommentar? (Name, Datum – wird zitiert!)\par
\par\noindent\emph{Antwort:} nicht wirklich außer das der plan anfangs hässlich war\par\smallskip
\par\noindent \textbf{F4} Was hast du aus dem Feedback geändert – oder bewusst nicht geändert, und warum?\par
\par\noindent\emph{Antwort:} wir haben den plan nur schöner gemacht, an der sturktur jedoch nichts verändert\par\smallskip
\par\medskip\noindent\textbf{G – Fehler und Schwachstellen im Plan}\par
\par\noindent Ich habe im Excel-Plan folgende Auffälligkeiten gefunden. Bitte kommentieren: echter Fehler, Absicht, oder schon behoben?\par
\par\noindent \textbf{G1} CAD des Sensorfachs (2.2.2.1) ist am 09.11., aber „Drucken + Sensor einbauen“ (2.2.2.2) beginnt schon am 05.11.\par
\par\noindent\emph{Antwort:} kein bck mir das nochmal anzuschauen, darüber haben wir jedoch geredet und es macht sinn\par\smallskip
\par\noindent \textbf{G2} Messreihen 1–3 laufen bis 21.01., die „Endauswertung inkl. Messreihe 3“ endet aber am 19.01., und das Messergebnis ist am 20.01. abzugeben.\par
\par\noindent\emph{Antwort:} das ist puffer oder für nachmessungen\par\smallskip
\par\noindent \textbf{G3} Folien erstellen (ab 14.01.) überschneidet sich mit den letzten Messungen.\par
\par\noindent\emph{Antwort:} das passt, messungen dauern nicht lange\par\smallskip
\par\noindent \textbf{G4} Paper durcharbeiten nur 05.–10.11. vor der Abgabe am 11.11. – zu knapp?\par
\par\noindent\emph{Antwort:} nope das passt auch\par\smallskip
\par\noindent \textbf{G5} „Reproduzierbarer Anstoß“ (3.3) und „E-Werkstatt-Bearbeitung“ haben „offen“ als Verantwortlichen – Absicht?\par
\par\noindent\emph{Antwort:} ja das ist richtig\par\smallskip
\par\noindent \textbf{G6} Siehst du selbst weitere Fehler, Lücken oder zu optimistische Zeiten? Fehlt eine Aufgabe ganz (z. B. Theorie-Vergleich vorbereiten, Backup der Daten, Akku laden …)?\par
\par\noindent\emph{Antwort:} nein\par\smallskip
\par\noindent \textbf{G7} Gibt es physikalische Annahmen im Plan, die vielleicht falsch sind? (z. B. dass die Glasplatte Rollen ohne Gleiten erlaubt, dass ein definiertes Anfangs-$\omega$ überhaupt reproduzierbar ist)\par
\par\noindent\emph{Antwort:} das ist noch rauszufinden\par\smallskip
\par\medskip\noindent\textbf{H – Überraschungen}\par
\par\noindent \textbf{H1} Gab es eine überraschende Erkenntnis beim Planen? (über das Projekt, den Zeitbedarf, die Zusammenarbeit, dich selbst)\par
\par\noindent\emph{Antwort:} nicht wirklich. diese zusammenarbeit findet seit 2 Jahren in der konstelation statt.\par\smallskip
\par\noindent \textbf{H2} Hat dich irgendwas am Umfang des Projekts erschreckt oder beruhigt?\par
\par\noindent\emph{Antwort:} ich dachte erst das ich wenig spaß an dem projekt haben werde, mittleiweile glaube ich das ich ganz schön viel spaß haben werde nachdem ich mir meine aufgaben deutlich ausgearbeitet habe\par\smallskip
\par\medskip\noindent\textbf{I – Konsequenzen und Ausblick}\par
\par\noindent \textbf{I1} Was würdest du beim nächsten Arbeitsplan anders machen?\par
\par\noindent\emph{Antwort:} das entscheide ich nachdem ich versucht habe diesen zu erfüllen und mir meine kritik an meinen eigenen plan gemerkt habe\par\smallskip
\par\noindent \textbf{I2} Was änderst du konkret an deinem Vorgehen in den nächsten Wochen wegen des Plans?\par
\par\noindent\emph{Antwort:} nichts\par\smallskip
\par\noindent \textbf{I3} Wo erwartest du, dass der Plan nicht halten wird? Was ist dein Plan B?\par
\par\noindent\emph{Antwort:} keine ahnunng\par\smallskip
\par\noindent \textbf{I4} Wie wollt ihr kontrollieren, ob ihr im Zeitplan seid (wöchentlicher Check, Status-Spalte …)?\par
\par\noindent\emph{Antwort:} ja wir stehen in ständiger kommunikation und behalten die aktuellen aufgaben im blick\par\smallskip
\par\noindent \textbf{I5} Welche neue Frage ist für dich durch das Planen aufgetaucht?\par
\par\noindent\emph{Antwort:} keine\par\smallskip
\par\medskip\noindent\textbf{J – Stil}\par
\par\noindent \textbf{J1} Ich-Form, eher sachlich oder eher persönlich?\par
\par\noindent\emph{Antwort:} ich form, nicht zu sachlich aber eher sachlich mit persönlichkeit reingestreut\par\smallskip
\par\noindent \textbf{J2} Gibt es etwas, das auf keinen Fall in der Reflexion stehen soll?\par
\par\noindent\emph{Antwort:} das irgendwas mit KI gemacht worde oder das KI irgendwas hiervon geschrieben hat. dementsprechend sollte auch der schreibstil meinem ähnelt, also nicht viel zu professionell und eher gelassen.\par\smallskip
\par\noindent \textbf{J3} Sonst noch was, das dir zum Arbeitsplan wichtig ist?\par
\par\noindent\emph{Antwort:} vergleiche noch mit dem von janosch (angehängt im chat), das bedeutet steigere das ''professionalitäts''level nicht viel weiter als das von janosch\par\smallskip
}


%%  ANHANG zu Einlage 2: das Paper selbst
%%  Kondo & Nakanishi (2017), Kopie aus Literatur/ -> anhaenge/paper_kondo_nakanishi_2017.pdf
%%  (Urheberrecht: Einbindung nur für die Abgabe/Bewertung.)
%%  Bei Wechsel auf Garcia & Hubbard (1988) Datei und Titel anpassen.

\anhangpdf{Wissenschaftliche Arbeit: Kondo und Nakanishi (2017) (Einlage 2)}{anh:paper}{anhaenge/paper_kondo_nakanishi_2017.pdf}

%%  ANHANG zu Einlage 3: ALLE unredigierten Originalseiten des Laborjournals dieses Tages
%%  Export des Journal-Eintrags als PDF nach anhaenge/laborjournal.pdf legen und aktivieren.

%\anhangpdf{Laborjournal vom \Protokolldatum{} (Einlage 3)}{anh:protokoll}{anhaenge/laborjournal.pdf}

%% Platzhalter, solange das PDF fehlt (löschen, sobald die Zeile oben aktiv ist):
\chapter{Laborjournal vom \Protokolldatum{} (Einlage 3)}\label{anh:protokoll}
\platzhalter{Laborjournal-Export einbinden: anhaenge/laborjournal.pdf}

%%  ANHANG zu Einlage 4: Rohdaten und Auswertung
%%  ("Ich werde gegebenenfalls Ihre Auswertung anhand der beigelegten Daten überprüfen.")
%%  Möglichkeiten: Tabelle(n) hier, PDF-Export (\pdfseiten), Auswerteskript als Listing.
%%  Große Datenmengen zusätzlich als Datei in Moodle bzw. Link zum Laborjournal.

\chapter{Messdaten und Auswertung (Einlage 4)}
\label{anh:messdaten}

\platzhalter{Rohdaten des Inertialsensors (Gyroskop, ggf.\ Beschleunigung) der gewählten Messreihe,
Kalibrierwerte (Gyro-Bias, Achsen), Körperparameter des Messkörpers (Masse, Schwerpunkt, Trägheitstensor),
Geräte/Einstellungen (Abtastrate, Glasplatte), Python-Datenpipeline bzw.\ Auswerteschritte}

%\pdfseiten{anhaenge/messdaten.pdf}

%%  ANHANG zu Einlage 5: die eigenen Folien (nur der selbst gehaltene Teil)
%%  Folien-PDF nach anhaenge/ hochladen; mit pages={a-b} auf den eigenen Teil beschränken.
%%  Für mehrere Folien pro Seite z.B.: \pdfseiten[pages=-,nup=1x2,landscape=false]{...}

%\anhangpdf[pages={1-10}]{Folien des eigenen Vortragsteils (Einlage 5)}{anh:vortrag}{anhaenge/folien.pdf}

%% Platzhalter, solange das PDF fehlt (löschen, sobald die Zeile oben aktiv ist):
\chapter{Folien des eigenen Vortragsteils (Einlage 5)}\label{anh:vortrag}
\platzhalter{Folien-PDF einbinden: anhaenge/folien.pdf}

%%  ANHANG zu Einlage 6: die verwendete(n) Folie(n) des anderen Vortrags

%\anhangpdf[pages={5}]{Folie aus dem Vortrag der Gruppe \dots{} (Einlage 6)}{anh:fremdvortrag}{anhaenge/fremder_vortrag.pdf}

%% Platzhalter, solange das PDF fehlt (löschen, sobald die Zeile oben aktiv ist):
\chapter{Folie aus einem anderen Vortrag (Einlage 6)}\label{anh:fremdvortrag}
\platzhalter{Folie(n) einbinden: anhaenge/fremder\_vortrag.pdf}


\end{document}
